\chapter*{解析几何-双曲线}
\subsection*{双曲线的定义}
我们把\textbf{平面上到两个不同定点 ${F}_{1}\text{ 、 }{F}_{2}$ 的距离之差的绝对值等于常数 (大于零且小于 $\left| {{F}_{1}{F}_{2}}\right|$ ) 的点的轨迹叫做双曲线}. 这两个定点 ${F}_{1}\text{ 、 }{F}_{2}$ 叫做双曲线的\textbf{焦点},两个焦点之间的距离 $\left| {{F}_{1}{F}_{2}}\right|$ 叫做双曲线的\textbf{焦距}.
\begin{example}
    将一条拉链先拉开一部分,并将拉开的一支的一端固定在 ${F}_{2}$ 处,在另一支上选择一点固定在 ${F}_{1}$ 处,使 $\left| {M{F}_{2}}\right|  - \left| {M{F}_{1}}\right|  = \left| {P{F}_{1}}\right|$ 为常数,在点 $M$ 处放上一支铅笔, 然后逐渐拉开拉链,让铅笔尖 $M$ 顺势移动,则笔尖 $M$ 画出的图形就是使 $\left| {M{F}_{2}}\right|  - \left| {M{F}_{1}}\right|$ 为常数 $\left| {P{F}_{1}}\right|$ 的点 $M$ 的轨迹的一部分.

    \twotu{\includegraphics[width=0.4\textwidth]{bodyxb/src/bo_d17r1jf7aajc7388ojlg_0_410_747_306_289_0.jpg}}{\includegraphics[width=0.4\textwidth]{bodyxb/src/bo_d17r1jf7aajc7388ojlg_0_896_784_308_255_0.jpg}}
\end{example}

\subsection*{双曲线的标准方程}
\begin{center}
\includegraphics[width=0.25\linewidth]{bodyxb/src/bo_d17r1jf7aajc7388ojlg_0_1185_1610_309_261_0.jpg}
\end{center}

设双曲线的焦距 $\left| {{F}_{1}{F}_{2}}\right|  = {2c}\left( {c > 0}\right)$ ,则点 ${F}_{1}\text{ 、 }{F}_{2}$ 的坐标分别为 $\left( {-c,0}\right) ,\left( {c,0}\right)$ .

设动点 $M\left( {x,y}\right)$ 是双曲线上任意一点,则点 $M$ 到两个定点 ${F}_{1}\text{ 、 }{F}_{2}$ 的距离之差的绝对值等于常数,设为 ${2a}\left( {0 < {2a} < {2c}}\right)$ ,即双曲线就是由下面的集合所描述的图形: $\left\{  {M\left| \right| \left| {M{F}_{1}}\right|  - \left| {M{F}_{2}}\right|  \mid   = {2a},{2c} > {2a} > 0}\right\}.$

所以有$\left| {M{F}_{1}}\right|  - \left| {M{F}_{2}}\right|  =  \pm  {2a}$，即$\sqrt{{\left( x + c\right) }^{2} + {y}^{2}} - \sqrt{{\left( x - c\right) }^{2} + {y}^{2}} =  \pm  {2a}$，即$\sqrt{{\left( x + c\right) }^{2} + {y}^{2}} = \sqrt{{\left( x - c\right) }^{2} + {y}^{2}} \pm  {2a}$

将上式两边平方,并整理后可得 ${cx} - {a}^{2} =  \pm  a\sqrt{{\left( x - c\right) }^{2} + {y}^{2}}$.

再将上式两边平方,并整理后可得 $\left( {{c}^{2} - {a}^{2}}\right) {x}^{2} - {a}^{2}{y}^{2} = {a}^{2}\left( {{c}^{2} - {a}^{2}}\right)$ ,因为 ${2c} > {2a} > 0$ ,即 $c >$  $a > 0$ ,所以 ${c}^{2} - {a}^{2} > 0$ ,令 ${b}^{2} = {c}^{2} - {a}^{2}$ ,且 $b > 0$，则得到双曲线的标准方程$\dfrac{x^2}{a^2} - \dfrac{y^2}{b^2} = 1\left( a > 0,b > 0\right).$

类似地,如果以焦点 ${F}_{1}\text{ 、 }{F}_{2}$ 所在直线为 $y$ 轴,线段 ${F}_{1}{F}_{2}$ 的垂直平分线为 $x$ 轴，建立平面直角坐标系如图所示,则可得到双曲线的方程为$\dfrac{{y}^{2}}{{a}^{2}} - \dfrac{{x}^{2}}{{b}^{2}} = 1\left( {a > 0,b > 0}\right)$

\begin{center}
\includegraphics[width=0.2\textwidth]{bodyxb/src/bo_d17r1jf7aajc7388ojlg_1_1178_939_284_357_0.jpg}
\end{center}


例题：
\begin{enumerate}
    \item 
    \begin{enumerate}
        \item 已知点 $M\left( {x,y}\right)$ 到点 ${F}_{1}\left( {-3,0}\right)$ 的距离减去它到点 ${F}_{2} (3,0)$的距离的差是 4,求点 $M$ 的轨迹方程
        
        \item 已知点 $M\left( {x,y}\right)$ 到点 ${F}_{1}\left( {-3,0}\right)$ 的距离与它到点 ${F}_{2}\left( {3,0}\right)$ 的距离的差的绝对值等于 6,求点 $M$ 的轨迹方程.
    \end{enumerate}
    \begin{jieda}
    \begin{enumerate}
        \item 依题意, $\left| {M{F}_{1}}\right|  - \left| {M{F}_{2}}\right|  = 4\left( {4 < \left| {{F}_{1}{F}_{2}}\right| }\right)$ ,所以点 $M$ 的轨迹是以 ${F}_{1}\text{ 、 }{F}_{2}$ 为焦点的双曲线的右支,且 $c = 3,{2a} = 4$ .
        
        因此 ${b}^{2} = {c}^{2} - {a}^{2} = 5$ ,所以点 $M$ 的轨迹方程为 $\dfrac{{x}^{2}}{4} - \dfrac{{y}^{2}}{5} = 1\left( {x \geq  2}\right)$ .
        \item 因为 $\left| \right| M{F}_{1}\left| -\right| M{F}_{2}\left| \right|  = 6$ ,而 $\left| {{F}_{1}{F}_{2}}\right|  = 6$ ,所以当 $\left| {M{F}_{1}}\right|  - \left| {M{F}_{2}}\right|  = 6$ 时,点 $M$ 的轨迹是在 $x$ 轴上以 ${F}_{2}$ 为端点向右的一条射线 $y = 0\left( {x \geq  3}\right)$ ;
        
        当 $\left| {M{F}_{1}}\right|  - \left| {M{F}_{2}}\right|  =  - 6$ 时,点 $M$ 的轨迹是在 $x$ 轴上以 ${F}_{1}$ 为端点向左的一条射线 $y =$  $0\left( {x \leq   - 3}\right)$ .
        
        所以点 $M$ 的轨迹方程为 $y = 0\left( {\left| x\right|  \geq  3}\right)$ .
    \end{enumerate}
    \end{jieda}
    \item 设 $P$ 是双曲线 $\dfrac{{x}^{2}}{16} - \dfrac{{y}^{2}}{20} = 1$ 上一点, ${F}_{1}$ , ${F}_{2}$ 分别是双曲线的左、右焦点,若 $\left| {P{F}_{1}}\right|  = 9$ ,则 $\left| {P{F}_{2}}\right|  =$ \blkda{17}
    
    \begin{jieda}
        根据双曲线的定义得 $||P{F}_{1}| - |P{F}_{2}||= 8$. 因为 $\left| {P{F}_{1}}\right|  = 9$，所以 $\left| {P{F}_{2}}\right|  = 1$ 或 17.
        
        又 $\left| {P{F}_{2}}\right|  \geq  c - a = 2$ ,故 $\left| {P{F}_{2}}\right| = 17$
    \end{jieda}
    \item 已知曲线 $C$ 的方程为 $\dfrac{{x}^{2}}{k + 1} + \dfrac{{y}^{2}}{5 - k} = 1\left( {k \in  \mathbf{R}}\right)$ ,若曲线 $C$ 是焦点在 $y$ 轴上的双曲线,则实数 $k$ 的取值范围是\blkda{$k <  - 1$}
    \begin{jieda}
        若曲线 $C$ 是焦点在 $y$ 轴上的双曲线，则 $\left\{  \begin{array}{l} k + 1 < 0, \\  5 - k > 0, \end{array}\right.$ 解得 $k <  - 1$ .
    \end{jieda}
    \item 方程 ${x}^{2} + \left( {\cos \theta }\right) {y}^{2} = 1,\theta  \in  \left( {0,\pi }\right)$ 表示的曲线不可能为\dotfill{(\kh{B})}
    \xx{两条直线}{圆}{椭圆}{双曲线}
    \begin{jieda}
        因为 $\theta  \in  \left( {0,\pi }\right)$ ,所以 $\cos \theta  \in  \left( {-1,1}\right)$
        \begin{dingautolist}{172}
            \item 当 $\cos \theta  \in  \left( {-1,0}\right)$ 时,方程 ${x}^{2} + \left( {\cos \theta }\right) {y}^{2} = 1$ 表示双曲线
            \item 当 $\cos \theta  = 0$ 时,方程 ${x}^{2} + \left( {\cos \theta }\right) {y}^{2} = 1$ 表示两条直线 $x =  \pm  1$
            \item 当 $\cos \theta  \in  \left( {0,1}\right)$ 时,方程 ${x}^{2} + \left( {\cos \theta }\right) {y}^{2} = 1$ 可化为 ${x}^{2} + \dfrac{{y}^{2}}{\dfrac{1}{\cos \theta }} = 1$
            
            因为 $\dfrac{1}{\cos \theta } > 1$ ,所以方程表示焦点在 $y$ 轴上的椭圆
        \end{dingautolist}
    \end{jieda}
    \item 
    \begin{enumerate}
        \item 求与双曲线 $C : \dfrac{{x}^{2}}{2} - {y}^{2} = 1$ 有公共的焦点,且过点 $P\left( {\sqrt{2},\sqrt{2}}\right)$ 的双曲线的标准方程
        \item 求经过两点 $P\left( {-\sqrt{3},2\sqrt{2}}\right)$ 、 $Q\left( {\sqrt{6}, - 2\sqrt{3}}\right)$ 的双曲线 $C$ 的标准方程.
    \end{enumerate}
    \begin{jieda}
        \begin{enumerate}
            \item 在双曲线 $C : \dfrac{{x}^{2}}{2} - {y}^{2} = 1$ 中, ${a}^{2} = 2,{b}^{2} = 1$ ,则 ${c}^{2} = {a}^{2} + {b}^{2} = 3$ ,所以双曲线 $C$ 的焦点坐标为 ${F}_{1}\left( {-\sqrt{3},0}\right) \text{ 、 }{F}_{2}\left( {\sqrt{3},0}\right)$ .
            
            设所求双曲线的方程为: $\dfrac{{x}^{2}}{{a}_{1}^{2}} - \dfrac{{y}^{2}}{{b}_{1}^{2}} = 1\left( {{a}_{1} > 0,{b}_{1} > 0}\right)$ ,则 ${c}_{1}^{2} = {a}_{1}^{2} + {b}_{1}^{2} = 3$.
            \begin{dingautolist}{172}
                \item 定义法
                
                由双曲线定义, 得${2a_1} = \left| \right| P{F}_{1}\left| -\right| P{F}_{2}\left| \right|  = \left| {\sqrt{{\left( \sqrt{2} + \sqrt{3}\right) }^{2} + 2} - \sqrt{{\left( \sqrt{2} - \sqrt{3}\right) }^{2} + 2}}\right|  = \left| {\left( {\sqrt{6} + 1}\right)  - \left( {\sqrt{6} - 1}\right) }\right|  = 2.$ 
                
                所以 $a_1 = 1,b_1 = \sqrt{2}$ . 

                \item 待定系数法

                代入$P$点坐标可得$\dfrac{2}{{a}_{1}^{2}} - \dfrac{2}{{b}_{1}^{2}} = 1 \left( {{a}_{1} > 0,{b}_{1} > 0}\right)$，故$\left\{\begin{array}{l}
                     {a}_{1}^{2} + {b}_{1}^{2} = 3  \\
                     \dfrac{1}{{a}_{1}^{2}} - \dfrac{1}{{b}_{1}^{2}} = \dfrac{1}{2}
                \end{array}\right.$，解得$\left\{\begin{array}{l}
                     a_1^2 = 1  \\
                     b_1^2 = 2
                \end{array}\right.$
            \end{dingautolist}
            因此，所求双曲线的标准方程为 ${x}^{2} - \dfrac{{y}^{2}}{2} = 1$ .
            
            \item 根据题意,可设双曲线 $C$ 的方程为 $A{x}^{2} - B{y}^{2} = 1\left( {{AB} > 0}\right)$ . 因为双曲线 $C$ 过 $P\left( {-\sqrt{3},2\sqrt{2}}\right) \text{ 、 }Q\left( {\sqrt{6}, - 2\sqrt{3}}\right)$ 两点,所以 $\left\{  \begin{array}{l} {3A} - {8B} = 1, \\  {6A} - {12B} = 1, \end{array}\right.$ 解得 $\left\{  \begin{array}{l} A =  - \dfrac{1}{3}, \\  B =  - \dfrac{1}{4}. \end{array}\right.$
            
            因此,双曲线的标准方程为 $\dfrac{{y}^{2}}{4} - \dfrac{{x}^{2}}{3} = 1$ .
        \end{enumerate}
    \end{jieda}
    \item 已知圆 ${C}_{1} : {\left( x + 3\right) }^{2} + {y}^{2} = 1,{C}_{2} : {\left( x - 3\right) }^{2} + {y}^{2} = 9$ ,动圆 $M$ 同时与圆 ${C}_{1}$ 和圆 ${C}_{2}$ 相外切,则动圆圆心 $M$ 的轨迹方程为\blkda[4]{${x}^{2} - \dfrac{{y}^{2}}{8} = 1\left( {x \leq   - 1}\right)$}

    \begin{jieda}
        设动圆 $M$ 的半径为 $r$ ,由动圆 $M$ 同时与圆 ${C}_{1}$ 和圆 ${C}_{2}$ 相外切，$\left| {M{C}_{1}}\right|  = 1 + r,\left| {M{C}_{2}}\right|  = 3 + r$，$\left| {M{C}_{2}}\right|  - \left| {M{C}_{1}}\right|  = 2 < 6,$
        
        所以动圆圆心 $M$ 的轨迹是以点 ${C}_{1}\left( {-3,0}\right)$ 和 ${C}_{2}\left( {3,0}\right)$ 为焦点的双曲线的左支，且 ${2a} = 2$ ,解得 $a = 1$ ,又 $c = 3$，则 ${b}^{2} = {c}^{2} - {a}^{2} = 8$ ,
        
        所以动圆圆心 $M$ 的轨迹方程为 ${x}^{2} - \dfrac{{y}^{2}}{8} = 1\left( {x \leq   - 1}\right)$ .
    \end{jieda}
    

    \item 设点 $P$ 是双曲线 $\dfrac{{x}^{2}}{{a}^{2}} - \dfrac{{y}^{2}}{{b}^{2}} = 1\left( {a > 0,b > 0}\right)$ 上任一点,且 $\angle {F}_{1}P{F}_{2} = {\theta }\left( {0 < \theta  < \pi}\right)$ , 求 $\bigtriangleup P{F}_{1}{F}_{2}$ 的面积.
    \begin{jieda}
        根据双曲线的定义有$||PF_1| - |PF_2|| = 2a$，平方有$|PF_1|^2+|PF_2|^2-2|PF_1||PF_2| = 4a^2$

        根据余弦定理有$|PF_1|^2 + |PF_2|^2 - 2|PF_1||PF_2|cos\theta = 4c^2$

        两式做差得到$2|PF_1||PF_2|(1-cos\theta) = 4b^2$

        根据半角公式变形有$\dfrac{1}{2}|PF_1||PF_2|sin\theta \cdot tan\dfrac{\theta}{2} = b^2$，即$S_{\triangle PF_1F_2} = \dfrac{b^2}{tan\dfrac{\theta}{2}} = b^2cot\dfrac{\theta}{2}$
    \end{jieda}
      \item 已知双曲线 $\dfrac{{x}^{2}}{4} - \dfrac{{y}^{2}}{5} = 1$ 的左焦点为 $F$ ,点 $P$ 在双曲线上且在 $x$ 轴上方,若线段 ${PF}$ 的中点在以 $O$ 为圆心, $\left| {OF}\right|$ 为半径的圆上,则直线 ${PF}$ 的斜率为\blkda[3]{$\dfrac{\sqrt{11}}{5}$ 或 $\sqrt{35}$}.
    \begin{jieda}
        由双曲线 $\dfrac{{x}^{2}}{4} - \dfrac{{y}^{2}}{5} = 1$ 可知 $a = 2,b = \sqrt{5},c = \sqrt{{a}^{2} + {b}^{2}} = 3$ ,
        
        设线段 ${PF}$ 的中点为 $M$ ,双曲线的右焦点为 ${A}$，则 $\left| {OM}\right|  = 3,\left| {AF}\right|  = 6$
        \begin{center}
            \includegraphics[width=0.3\textwidth]{bodyxb/src/bo_d3eri7s601uc738iqr70_2_240_1616_444_409_0.jpg}
        \end{center}
        \begin{dingautolist}{172}
            \item 若点 $P$ 在第一象限,则 $\left| {PA}\right|  = 2\left| {OM}\right|  = 6,\left| {PF}\right|  = \left| {PA}\right|  + 4 = {10}$，可得 $\cos \angle {PFA} = \dfrac{{\left| PA\right| }^{2} + {\left| AF\right| }^{2} - {\left| AP\right| }^{2}}{2\left| {PA}\right|  \cdot  \left| {AF}\right| } = \dfrac{{10}^{2} + {6}^{2} - {6}^{2}}{2 \times  {10} \times  6} = \dfrac{5}{6}$，则 $\sin \angle {PFA} = \sqrt{1 - {\cos }^{2}\angle {PFA}} = \dfrac{\sqrt{11}}{6}$
            
            可得 $\tan \angle {PFA} = \dfrac{\sin \angle {PFA}}{\cos \angle {PFA}} = \dfrac{\sqrt{11}}{5}$ ,所以直线 ${PF}$ 的斜率为 $\dfrac{\sqrt{11}}{5}$ .
            
            \item 若点 $P$ 在第二象限，则$\left| {PA}\right|  = 2\left| {OM}\right|  = 6,\left| {PF}\right|  = \left| {PA}\right|  - 4 = 2$

            ，可得 $\cos \angle {PFA} = \dfrac{{\left| PA\right| }^{2} + {\left| AF\right| }^{2} - {\left| AP\right| }^{2}}{2\left| {PA}\right|  \cdot  \left| {AF}\right| } = \dfrac{2^{2} + {6}^{2} - {6}^{2}}{2 \times  2 \times  6} = \dfrac{1}{6}$，$\sin \angle {PFA} = \sqrt{1 - {\cos }^{2}\angle {PFA}} = \dfrac{\sqrt{35}}{6}$

            可得 $\tan \angle {PFA} = \dfrac{\sin \angle {PFA}}{\cos \angle {PFA}} = \sqrt{35}$ ,所以直线 ${PF}$ 的斜率为 $\sqrt{35}$ .
        
        \end{dingautolist}
        \ps 直接求解$\left\{  \begin{array}{l} \dfrac{{\left( s - 3\right) }^{2}}{4} + \dfrac{81}{4} = 9 \\  \dfrac{{s}^{2}}{4} - \dfrac{81}{5} = 1  \\  t > 0 \end{array}\right.$即可得$P$在一二两个象限的坐标$(s, t)$.
    \end{jieda}
\end{enumerate}


\subsection*{双曲线的性质}

下面根据双曲线 $C$ 的标准方程$\dfrac{{x}^{2}}{{a}^{2}} - \dfrac{{y}^{2}}{{b}^{2}} = 1\left( {a > 0,b > 0}\right)$来研究它的性质.

\subsubsection*{对称性}

与研究椭圆的对称性类似,可知双曲线既是轴对称图形,也是中心对称图形,它有两条互相垂直的对称轴,其唯一的对称中心叫做双曲线的中心. 在双曲线的标准方程的条件下,双曲线的对称轴就是坐标轴,双曲线的中心就是坐标原点.

\subsubsection*{顶点}

双曲线与它的对称轴的交点叫做双曲线的顶点.在方程$\dfrac{{x}^{2}}{{a}^{2}} - \dfrac{{y}^{2}}{{b}^{2}} = 1\left( {a > 0,b > 0}\right)$中
\begin{dingautolist}{172}
    \item 令 $y = 0$，得 $x =  \pm  a$ ,因此 ${A}_{1}\left( -a,0\right)$ 、 ${A}_{2}\left( a,0\right)$ 是双曲线的两个顶点. 
    \item 令 $x = 0$，得 ${y}^{2} =  - {b}^{2}$，此时$y$ 没有实数根,因此双曲线与 $y$ 轴没有交点，但我们也把 ${B}_{1}\left( {0, - b}\right) ,{B}_{2}\left( {0,b}\right)$ 画在 $y$ 轴上,如图所示.
\end{dingautolist}
\begin{center}
\includegraphics[width=0.2\textwidth]{bodyxb/src/bo_d17r1jf7aajc7388ojlg_5_1205_1034_270_314_0.jpg}
\end{center}

线段 ${A}_{1}{A}_{2}$ 和 ${B}_{1}{B}_{2}$ 分别在双曲线的两条对称轴上,我们把线段 ${A}_{1}{A}_{2}$ 叫做双曲线的实轴,实轴的长等于 $2a$，$a$ 为双曲线的实半轴长; 

线段 ${B}_{1}{B}_{2}$ 叫做双曲线的虚轴,虚轴的长等于 ${2b},b$ 为双曲线的虚半轴长. 

双曲线的焦距为 ${2c}$ ,半焦距为 $c$ ,且满足 ${c}^{2} = {a}^{2} + {b}^{2}$ . 显然，双曲线的两个焦点都在它的实轴上.

特别地，我们把实轴和虚轴等长的双曲线叫做\textbf{等轴双曲线},即在标准方程中 $a = b$ ；我们把实轴和虚轴对换的两个双曲线叫做一对\textbf{共轭双曲线},即方程$\dfrac{{x}^{2}}{{a}^{2}} - \dfrac{{y}^{2}}{{b}^{2}} = 1\left( {a > 0,b > 0}\right)$与方程 $\dfrac{{x}^{2}}{{a}^{2}} - \dfrac{{y}^{2}}{{b}^{2}} =  - 1\left( {a > 0,b > 0}\right)$ 代表一对共轭双曲线.

\subsubsection*{范围}
由方程$\dfrac{{x}^{2}}{{a}^{2}} - \dfrac{{y}^{2}}{{b}^{2}} = 1\left( {a > 0,b > 0}\right)$可知 $\dfrac{{x}^{2}}{{a}^{2}} - 1 = \dfrac{{y}^{2}}{{b}^{2}} \geq  0$ ,所以 $\dfrac{{x}^{2}}{{a}^{2}} \geq  1$ ,即 $\left| x\right|  \geq  a$ ,且随着 $\left| x\right|$ 不断增大, $\left| y\right|$ 也相应地不断增大. 所以,双曲线 $\dfrac{{x}^{2}}{{a}^{2}} - \dfrac{{y}^{2}}{{b}^{2}} = 1(a > 0$ , $b > 0$ ) 在不等式 $x \geq  a$ 或 $x \leq   - a$ 所表示的区域内,且由两支组成:一支双曲线在直线 $x =  - a$ 的左侧,向左上、左下两方无限延伸; 另一支双曲线在直线 $x = a$ 的右侧,向右上、右下两方无限延伸.

\begin{center}
\includegraphics[width=0.2\textwidth]{bodyxb/src/bo_d17r1jf7aajc7388ojlg_5_1180_1728_303_298_0.jpg}
\end{center}

\subsubsection*{渐近线}

过双曲线的实轴顶点 ${A}_{1}\left( {-a,0}\right) \text{ 、 }{A}_{2}\left( {a,0}\right)$ 作 $y$ 轴的平行线 $x =  \pm  a$ ,过双曲线的虚轴顶点 ${B}_{1}\left( {0, - b}\right) \text{ 、 }{B}_{2}\left( {0,b}\right)$ 作 $x$ 轴的平行线 $y =  \pm  b$ ,它们围成一个矩形,如图所示. 

\begin{center}
\includegraphics[width=0.2\textwidth]{bodyxb/src/bo_d17r1jf7aajc7388ojlg_6_1157_270_304_303_0.jpg}
\end{center}

这个矩形的两条对角线所在的直线的方程是 $y =  \pm  \dfrac{b}{a}x$ ,当双曲线的两支向外无限延伸时,双曲线与这两条直线无限趋近,但永不相交. 所以，称直线 $y =  \pm  \dfrac{b}{a}x$ 为双曲线 $\dfrac{{x}^{2}}{{a}^{2}} - \dfrac{{y}^{2}}{{b}^{2}} = 1\left( {a > 0\text{ 、 }b > 0}\right)$ 的两条渐近线，而双曲线位于以两条渐近线为边界的左、右平面区域内.

证明如下：

双曲线在第一象限内部分的方程为 $y = \dfrac{b}{a}\sqrt{{x}^{2} - {a}^{2}}\left( {x > a}\right)$ . 如图所示，设点 $M\left( {{x}_{0},{y}_{0}}\right)$ 是第一象限内双曲线上面的任意一点,且 ${x}_{0} > a > 0$ ; 点 $P\left( {{x}_{0},{y}_{1}}\right)$ 是直线 $y = \dfrac{b}{a}x$ 上与 $M$ 有相同横坐标的点,则 ${y}_{1} = \dfrac{b}{a}{x}_{0}$ .

\begin{center}
\includegraphics[width=0.3\textwidth]{bodyxb/src/bo_d17r1jf7aajc7388ojlg_6_1060_879_413_292_0.jpg}
\end{center}

因为 ${y}_{0} = \dfrac{b}{a}\sqrt{{x}_{0}^{2} - {a}^{2}} = \dfrac{b}{a}{x}_{0}\sqrt{1 - {\left( \dfrac{a}{{x}_{0}}\right) }^{2}} < \dfrac{b}{a}{x}_{0} = {y}_{1}$ , 所以$\left| {MP}\right|  = {y}_{1} - {y}_{0} = \dfrac{b}{a}\left( {{x}_{0} - \sqrt{{x}_{0}^{2} - {a}^{2}}}\right)  = \dfrac{b}{a} \cdot  \dfrac{\left( {{x}_{0} - \sqrt{{x}_{0}^{2} - {a}^{2}}}\right) \left( {{x}_{0} + \sqrt{{x}_{0}^{2} - {a}^{2}}}\right) }{{x}_{0} + \sqrt{{x}_{0}^{2} - {a}^{2}}} = \dfrac{ab}{{x}_{0} + \sqrt{{x}_{0}^{2} - {a}^{2}}}.
$

当 ${x}_{0} \Longrightarrow     + \infty$ 时, $\left| {MP}\right|  \Longrightarrow    0$ ,即双曲线的右支向右上方无限延伸时,它总在渐近线 $y = \dfrac{b}{a}x$ 的下方,与渐近线 $y = \dfrac{b}{a}x$ 无限趋近,但永不相交.

根据双曲线的对称性可知，双曲线 $\dfrac{{x}^{2}}{{a}^{2}} - \dfrac{{y}^{2}}{{b}^{2}} = 1\left( {a > 0,b > 0}\right)$ 和两条渐近线在其他象限也有类似性质. 可以证明, 对于焦点在 $y$ 轴上的双曲线 $\dfrac{{y}^{2}}{{a}^{2}} - \dfrac{{x}^{2}}{{b}^{2}} = 1$ ,其渐近线相应的为 $x = \pm  \dfrac{b}{a}y$，即 $y =  \pm  \dfrac{a}{b}x$.

结论：
\begin{enumerate}
    \item 双曲线$Ax^2+By^2 = 1 (AB < 0)$的渐近线为$Ax^2+By^2 = 0 (AB < 0)$
    \item 等轴双曲线（$x^2-y^2 = A$）的渐近线斜率为$\pm 1$，共轭双曲线（$\dfrac{{x}^{2}}{{a}^{2}} - \dfrac{{y}^{2}}{{b}^{2}} = 1$与$\dfrac{{y}^{2}}{{b}^{2}} - \dfrac{{x}^{2}}{{a}^{2}} = 1$）的渐近线相同，$\dfrac{{x}^{2}}{{a}^{2}} - \dfrac{{y}^{2}}{{b}^{2}} = 1$与$\dfrac{{y}^{2}}{{a}^{2}} - \dfrac{{x}^{2}}{{b}^{2}} =  1$）的渐近线斜率互为倒数.
    \item 与双曲线 $Ax^2+By^2 = 1 (AB < 0)$ 有共同渐近线的方程可表示为 $Ax^2+By^2 = t (AB < 0) \left( {t \neq  0}\right)$.
    \item 渐近线与实轴夹角正切值为$\dfrac{b}{a}$，正弦值为$\dfrac{b}{c}$，余弦值为$\dfrac{a}{c}$，因此焦点到渐近线的距离为$b$.
    \item 计算两渐近线之间的夹角：
    
    $\dfrac{{x}^{2}}{{a}^{2}} - \dfrac{{y}^{2}}{{b}^{2}} = 1\left( {a > 0,b > 0}\right)$的两条渐近线为$y = \pm \dfrac{b}{a}x$，写成一般式为$ay \pm bx = 0$，设两渐近线的夹角为$\theta$
    \begin{dingautolist}{172}
        \item 计算$tan\theta$
    
        两渐近线的斜率为$k_1 = \dfrac{b}{a}, k_2 = -\dfrac{b}{a}$
        
        $tan\theta = \left|\dfrac{k_1 - k_2}{1 + k_1k_2}\right| = \left|\dfrac{\dfrac{2b}{a}}{1 - \dfrac{b^2}{a^2}}\right| = \left|\dfrac{2ab}{a^2-b^2}\right|$
        \item 计算$cos\theta$
    
        两渐近线的方向向量为$\vv{d_1} = (a, b), \vv{d_2} = (-a, b)$
    
        $cos\theta = \dfrac{|b^2-a^2|}{a^2+b^2}$
    \end{dingautolist}
\end{enumerate}




\subsubsection*{离心率}

双曲线的焦距与实轴长的比 $\dfrac{c}{a}$叫做双曲线的离心率，记作 $e = \dfrac{c}{a}$，因为 $c > a > 0$ ,所以双曲线的离心率 $e > 1$ . 

由等式 ${c}^{2} - {a}^{2} = {b}^{2}$ ,可得$\dfrac{b}{a} = \dfrac{\sqrt{{c}^{2} - {a}^{2}}}{a} = \sqrt{{\left( \dfrac{c}{a}\right) }^{2} - 1} = \sqrt{{e}^{2} - 1}.$

结论：
\begin{dingautolist}{172}
    \item 双曲线的离心率越大，它的开口就越开阔；反之，离心率越小，双曲线的开口就越狭窄
    \item $\dfrac{{x}^{2}}{{a}^{2}} - \dfrac{{y}^{2}}{{b}^{2}} = 1\left( {a > 0,b > 0}\right)$的渐近线的斜率$k$与离心率满足$k^2+1 = e^2$
\end{dingautolist}

\begin{center}
\adjustbox{width=\textwidth}{
\begin{tabular}{|c|c|c|c|}
\hline
\multicolumn{2}{|c|}{标准方程} & $\dfrac{{x}^{2}}{{a}^{2}} - \dfrac{{y}^{2}}{{b}^{2}} = 1\left( {a > 0,b > 0}\right)$ & $\dfrac{{y}^{2}}{{a}^{2}} - \dfrac{{x}^{2}}{{b}^{2}} = 1\left( {a > 0,b > 0}\right)$ \\
\cline{1-4}
\multicolumn{2}{|c|}{图形} &  \includegraphics[width=0.2\textwidth]{bodyxb/src/bo_d1chk83ef24c73d2s6gg_0_624_857_239_159_0.jpg}  &  \includegraphics[width=0.2\textwidth]{bodyxb/src/bo_d1chk83ef24c73d2s6gg_0_1042_851_203_176_0.jpg}  \\
\cline{1-4}
\multirow{8}{*}{性质} & 焦点 & ${F}_{1}\left( {-c,0}\right) ,{F}_{2}\left( {c,0}\right)$ & ${F}_{1}\left( {0, - c}\right) ,{F}_{2}\left( {0,c}\right)$ \\
\cline{2-4}
 & 焦距 & \multicolumn{2}{|c|}{$\left| {{F}_{1}{F}_{2}}\right|  = {2c}$} \\
\cline{2-4}
 & 范围 & $x \leq   - a$ 或 $x \geq  a,y \in  \mathbf{R}$ & $y \leq   - a$ 或 $y \geq  a,x \in  \mathbf{R}$ \\
\cline{2-4}
 & 对称性 & \multicolumn{2}{|c|}{对称轴:坐标轴；对称中心:原点} \\
\cline{2-4}
 & 顶点 & ${A}_{1}\left( {-a,0}\right) ,{A}_{2}\left( {a,0}\right)$ & ${A}_{1}\left( {0, - a}\right) ,{A}_{2}\left( {0,a}\right)$ \\
\cline{2-4}
 & 轴 & \multicolumn{2}{|c|}{实轴:线段 ${A}_{1}{A}_{2}$ ,长: ${2a}$ ；虚轴:线段 ${B}_{1}{B}_{2}$ ,长: ${2b}$ ,实半 轴长: $a$ ,虚半轴长: $b$} \\
\cline{2-4}
 & 渐近线 & $y =  \pm  \dfrac{b}{a}x$ & $y =  \pm  \dfrac{a}{b}x$ \\
\cline{2-4}
 & 离心率 & \multicolumn{2}{|c|}{$e = \dfrac{c}{a} \in  \left( {1, + \infty }\right)$} \\
\cline{1-4}
\multicolumn{2}{|c|}{$a,b,c$ 的关系} & \multicolumn{2}{|c|}{${c}^{2} = {a}^{2} + {b}^{2}\left( {c > a > 0,c > b > 0}\right)$} \\
\cline{1-4}
\hline
\end{tabular}
}
\end{center}

例题：
\begin{enumerate}
    \item 下列结论中正确的是\blkda{\ding{174}\ding{175}}
    \begin{dingautolist}{172}
        \item 平面内到点 ${F}_{1}\left( {0,4}\right) ,{F}_{2}\left( {0, - 4}\right)$ 的距离之差的绝对值等于 8 的点的轨迹是双曲线.
        \item 方程 $\dfrac{{x}^{2}}{m} - \dfrac{{y}^{2}}{n} = 1\left( {{mn} > 0}\right)$ 表示焦点在 $x$ 轴上的双曲线. 
        \item 双曲线 $\dfrac{{x}^{2}}{{m}^{2}} - \dfrac{{y}^{2}}{{n}^{2}} = 1\left( {m > 0,n > 0}\right)$ 的渐近线方程是 $\dfrac{x}{m} \pm  \dfrac{y}{n} = 0.$ 
        \item 等轴双曲线的渐近线互相垂直,离心率等于 $\sqrt{2}$
    \end{dingautolist}
    
    \item 求下列双曲线的实轴和虚轴的长、离心率、焦点和顶点坐标、渐近线方程:
    \begin{enumerate}
        \item ${x}^{2} - 7{y}^{2} = 7$
        \item ${x}^{2} - 2{y}^{2} =  -8$
    \end{enumerate}
    \begin{jieda}
        \begin{enumerate}
            \item $\dfrac{x^2}{7} - {y}^{2} = 1$，$a = \sqrt{7}, b = 1, c = 2\sqrt{2}$，实轴长$2\sqrt{7}$，虚轴长$2$，离心率$\dfrac{2\sqrt{14}}{7}$，焦点$(-2\sqrt{2}, 0)$和$(2\sqrt{2}, 0)$，顶点$(-\sqrt{7}, 0)$和$(\sqrt{7}, 0)$，渐近线方程为$y = \dfrac{\sqrt{7}}{7}x$和$y = -\dfrac{\sqrt{7}}{7}x$
            \item $\dfrac{y^2}{4} - \dfrac{x^2}{8} = 1$，$a = 2, b = 2\sqrt{2}, c = 2\sqrt{3}$，实轴长$4$，虚轴长$4\sqrt{2}$，离心率$\sqrt{3}$，焦点$(0, 2\sqrt{3})$和$(0, -2\sqrt{3})$，顶点$(0, 2)$和$(0, -2)$，渐近线方程$y = \dfrac{\sqrt{2}}{2}x$或$y = -\dfrac{\sqrt{2}}{2}x$
        \end{enumerate}
    \end{jieda}
    \item 已知直线 $y = x$ 与双曲线 $\dfrac{{x}^{2}}{{a}^{2}} - \dfrac{{y}^{2}}{{b}^{2}} = 1\left( {a > 0,b > 0}\right)$ 无公共点,则双曲线的离心率的取值范围是\blkda{$(1, \sqrt{2}]$}
    \begin{jieda}
        双曲线的一条渐近线为 $y = \dfrac{b}{a}x$，因为直线 $y = x$ 与双曲线无公共点，故有 $\dfrac{b}{a} \leq  1$，即 $\dfrac{{b}^{2}}{{a}^{2}} = \dfrac{{c}^{2} - {a}^{2}}{{a}^{2}} = {e}^{2} - 1 \leq  1$
        
        所以 ${e}^{2} \leq  2$ ,所以 $1 < e \leq  \sqrt{2}$ .
    \end{jieda}
    \item 双曲线 $2{y}^{2} - {x}^{2} = 1$ 的渐近线方程是\dotfill{(\kh{C})}
    \xx{$y =  \pm  \dfrac{1}{2}x$}{$y =  \pm  {2x}$}{$y =  \pm  \dfrac{\sqrt{2}}{2}x$}{$y =  \pm  \sqrt{2}x$}

    \begin{jieda}
        双曲线 $2{y}^{2} - {x}^{2} = 1$ 的渐近线方程是 $2{y}^{2} - {x}^{2} = 0$，即$y =  \pm  \dfrac{\sqrt{2}}{2}x$.
    \end{jieda}


    \item 求满足下列条件的双曲线的标准方程:
    \begin{enumerate}
        \item 与双曲线 $\dfrac{{x}^{2}}{9} - \dfrac{{y}^{2}}{16} = 1$ 有共同渐近线,且过点 $\left( {-3,2\sqrt{3}}\right)$
        \item 与双曲线 $\dfrac{{x}^{2}}{16} - \dfrac{{y}^{2}}{4} = 1$ 有公共焦点,且过点 $\left( {3\sqrt{2},2}\right)$ .
    \end{enumerate}
    \begin{jieda}
        \begin{enumerate}
            \item 设所求双曲线方程为 $\dfrac{{x}^{2}}{9} - \dfrac{{y}^{2}}{16} = \lambda \left( {\lambda  \neq  0}\right)$ . 因为双曲线过点 $\left( {-3,2\sqrt{3}}\right)$ ,所以 $\dfrac{9}{9} - \dfrac{12}{16} = \lambda$ ,解得 $\lambda  = \dfrac{1}{4}$ ,所以有 $\dfrac{{x}^{2}}{9} - \dfrac{{y}^{2}}{16} = \dfrac{1}{4}$ . 故所求双曲线的标准方程为 $\dfrac{{x}^{2}}{\dfrac{9}{4}} - \dfrac{{y}^{2}}{4} = 1$ .
            \item 设所求双曲线方程为 $\dfrac{{x}^{2}}{{16} - k} - \dfrac{{y}^{2}}{4 + k} = 1\left( {-4 < k < {16}}\right)$ . 因为双曲线过点 $\left( {3\sqrt{2},2}\right)$ , 所以 $\dfrac{18}{{16} - k} - \dfrac{4}{4 + k} = 1$ ,得 ${k}^{2} + {10k} - {56} = 0$ ,解得 $k = 4$ 或 $k =  - {14}$ (舍).
            
            故所求双曲线的标准方程为 $\dfrac{{x}^{2}}{12} - \dfrac{{y}^{2}}{8} = 1$ .
        \end{enumerate}
    \end{jieda}
    
    \item 在平面直角坐标系中,已知双曲线 $\dfrac{{x}^{2}}{{a}^{2}} - \dfrac{{y}^{2}}{{b}^{2}} = 1\left( {a > 0,b > 0}\right)$ 的右焦点为 $F$ ,点 $A$ 在双曲线的渐近线上, $\bigtriangleup {OAF}$ 是边长为 2 的等边三角形,则双曲线的标准方程为\blkda[3]{${x}^{2} - \dfrac{{y}^{2}}{3} = 1$}
    
    \begin{jieda}
        由方程 $\dfrac{{x}^{2}}{{a}^{2}} - \dfrac{{y}^{2}}{{b}^{2}} = 1$，得双曲线的渐近线方程为 $y =  \pm  \dfrac{b}{a}x$
        
        不妨设 $A$ 在直线 $y = \dfrac{b}{a}x$ 上，由 $\bigtriangleup {OAF}$ 是边长为 2 的等边三角形可得 $c = 2$ ,直线 $y = \dfrac{b}{a}x$ 的倾斜角为 ${60}^{ \circ  }$，即 $\dfrac{b}{a} = \sqrt{3}$ ,
        
        联立 $\left\{  {\begin{array}{l} b = \sqrt{3}a, \\  {a}^{2} + {b}^{2} = {c}^{2} = 4, \end{array}\text{ 可得 }\left\{  \begin{array}{l} b = \sqrt{3}, \\  a = 1, \end{array}\right. }\right.$
        
        故双曲线的标准方程为 ${x}^{2} - \dfrac{{y}^{2}}{3} = 1$ .
    \end{jieda}
    \item 
    \begin{enumerate}
        \item 已知双曲线 $C : \dfrac{{x}^{2}}{{a}^{2}} - \dfrac{{y}^{2}}{{b}^{2}} = 1\left( {a > 0,b > 0}\right)$ 的离心率 $e = 2$ ,则该双曲线 $C$ 的渐近线方程为\blkda[3]{$y =  \pm  \sqrt{3}x$}.
        \item 已知双曲线 $C : \dfrac{{y}^{2}}{{a}^{2}} - \dfrac{{x}^{2}}{{b}^{2}} = 1\left( {a > 0,b > 0}\right)$ 的离心率 $e = 2$ ,则该双曲线 $C$ 的渐近线方程为\blkda[3]{$y =  \pm  \dfrac{\sqrt{3}}{3}x$}.
    \end{enumerate}
    

    \begin{jieda}
        \begin{enumerate}
            \item 双曲线的渐近线为$\dfrac{{x}^{2}}{{a}^{2}} - \dfrac{{y}^{2}}{{b}^{2}} = 0$，故其斜率平方$k^2 = \dfrac{b^2}{a^2} = \dfrac{c^2-a^2}{a^2} = e^2 - 1 = 3$
        
            所以该双曲线的渐近线方程为 $y =  \pm  \dfrac{b}{a}x =  \pm  \sqrt{3}x$ .
            \item 双曲线的渐近线为$\dfrac{{y}^{2}}{{a}^{2}} - \dfrac{{x}^{2}}{{b}^{2}} = 0$，故其斜率平方$k^2 = \dfrac{a^2}{b^2} = \dfrac{a^2}{c^2-a^2} = \dfrac{1}{e^2 - 1} = \dfrac{1}{3}$
        
            所以该双曲线的渐近线方程为 $y =  \pm  \dfrac{a}{b}x =  \pm  \dfrac{\sqrt{3}}{3}x$ .
        \end{enumerate}
        
    \end{jieda}


    \item 已知双曲线 $C : \dfrac{{x}^{2}}{{a}^{2}} - \dfrac{{y}^{2}}{{b}^{2}} = 1\left( {a > 0,b > 0}\right)$ 的左、右焦点分别 ${F}_{1},{F}_{2}$，${A}$ 是 $C$ 上一点(在第一象限),直线 $A{F}_{2}$ 与 $y$ 轴交于点 $B$ ,若 $A{F}_{1} \bot BF_1$且$3\left| {A{F}_{2}}\right|  = 2\left| {{F}_{2}B}\right|$ ,则 $C$ 的渐近线方程为\blkda[3]{$y =  \pm  \dfrac{2\sqrt{5}}{5}x$}
    \begin{jieda}
        示意图如下：
        \begin{center}
            \includegraphics[width=0.25\linewidth]{bodyxb//src/jianjinxian2.png}
        \end{center}
    
        设 $\left| {B{F}_{1}}\right|  = m$ ,则 $\left| {B{F}_{2}}\right|  = m$ ,因为 $3\left| {A{F}_{2}}\right|  = 2\left| {{F}_{2}B}\right|$ ,
        所以 $\left| {A{F}_{2}}\right|  = \dfrac{2}{3}m$ ,根据双曲线的定义： $\left| {A{F}_{1}}\right|  = \dfrac{2}{3}m + {2a}$

        因为 $A{F}_{1} \bot  B{F}_{1}$ ,由勾股定理得： ${\left( \dfrac{2}{3}m + 2a\right) }^{2} + {m^2} = {\left( \dfrac{5}{3}m\right) }^{2} \Longrightarrow    \left( {m - {3a}}\right) \left( {m + a}\right)  = 0$ ,
        所以 $m = {3a}$ ,则 $\left| {B{F}_{1}}\right| = \left| {B{F}_{2}}\right| = {3a}$ , $\left| {A{F}_{2}}\right| = {2a}$ , $\left| {A{F}_{1}}\right| = {4a}$ 
        
        在 ${\triangle A}{F}_{1}{F}_{2}$ 中, $\cos \angle A{F}_{1}{F}_{2} = \dfrac{{\left| A{F}_{1}\right| }^{2} + {\left| {F}_{1}{F}_{2}\right| }^{2} - {\left| A{F}_{2}\right| }^{2}}{2\left| {A{F}_{1}}\right|  \cdot  \left| {{F}_{1}{F}_{2}}\right| } = \dfrac{3{a}^{2} + {c}^{2}}{4ac}$ .
        
        在 $\bigtriangleup BF_1F_2$ 中, $\cos \angle B{F}_{1}F_2 = \dfrac{{\left| B{F}_{1}\right| }^{2} + {\left| {F}_{1}{F}_{2}\right| }^{2} - {\left| B{F}_{2}\right| }^{2}}{2\left| {B{F}_{1}}\right|  \cdot  \left| {{F}_{1}{F}_{2}}\right| } = \dfrac{c}{3a}$ .

        因为$\angle{A{F}_{1}F_2} + \angle{B{F}_{1}F_2} = \ang{90}$，所以$\left(\dfrac{3{a}^{2} + {c}^{2}}{4ac}\right)^2 + \left(\dfrac{c}{3a}\right)^2 = 1$

        整理有$25e^4-90e^2+81 = 0$，解得$e^2 = \dfrac{9}{5}$，$\dfrac{b}{a} = \sqrt{e^2-1} = \dfrac{2\sqrt{5}}{5}$，
        所以双曲线渐近线的方程为: $y =  \pm  \dfrac{2\sqrt{5}}{5}x$ .
    \end{jieda}

    

    
    \item 双曲线 $C : \dfrac{{x}^{2}}{{a}^{2}} - \dfrac{{y}^{2}}{{b}^{2}} = 1\left( {a > 0,b > 0}\right)$ 上任意一点到其渐近线 ${l}_{1} : y = \dfrac{b}{a}x$ 和 ${l}_{2}$ : $y =  - \dfrac{b}{a}x$ 的距离之积是常数.
    \begin{jieda}
        设双曲线上点$P(x_0, y_0)$，则其到两条渐近线的距离分别是$\dfrac{|ay_0-bx_0|}{\sqrt{a^2+b^2}}$和$\dfrac{|ay_0+bx_0|}{\sqrt{a^2+b^2}}$，其乘积为$\dfrac{|a^2y_0^2 - b^2x_0^2|}{a^2+b^2} = \dfrac{a^2b^2}{a^2+b^2}$
    \end{jieda}
    
    \item 某电厂冷却塔的外形是由双曲线 ${C}_{1}$ 的一部分绕其虚轴所在的直线旋转所形成的曲面,如图所示,已知它的最小半径为 ${20}{m}$ ,上口半径为 ${10}\sqrt{5}{m}$，下口半径为 ${20}\sqrt{2}{m}$ ,高为 ${60}{m}$，以冷却塔的轴截面的最窄处所在的直线为 $x$ 轴且与曲面交于点 $A{A}^{\prime }$ ,以线段 $A{A}^{\prime }$ 的垂直平分线为 $y$ 轴建立平面直角坐标系.
    \begin{center}
    \includegraphics[width=0.5\textwidth]{bodyxb/src/bo_d17r1jf7aajc7388ojlg_8_433_1272_768_453_0.jpg}
    \end{center}
    \begin{enumerate}
        \item 求此双曲线 ${C}_{1}$ 的方程
        \item 求双曲线 ${C}_{1}$ 的共轭双曲线 ${C}_{2}$ 的方程
        \item 设(2)中的双曲线 ${C}_{1}$ 、 ${C}_{2}$ 的离心率分别为 ${e}_{1}$ 、 ${e}_{2}$ ,请写出 ${e}_{1}$ 与 ${e}_{2}$ 满足的一个关系式, 并证明之.
    \end{enumerate}
    \begin{jieda}
        \begin{enumerate}
            \item 设双曲线的方程为 $\dfrac{{x}^{2}}{{a}^{2}} - \dfrac{{y}^{2}}{{b}^{2}} = 1\left( {a > 0,b > 0}\right)$，由题意知，$\left| {A{A}^{\prime }}\right|  = {2a} = {40},a = {20},\left| {C{C}^{\prime }}\right|  = {20}\sqrt{5},\left| {B{B}^{\prime }}\right|  = {40}\sqrt{2}$
    
            所以，设点 $C\left( {{10}\sqrt{5},t}\right) \left( {t > 0}\right)$，点 $B\left( {{20}\sqrt{2}, - {60} + t}\right)$，
            
            则$\left\{  \begin{array}{l} \dfrac{{\left( {10}\sqrt{5}\right) }^{2}}{{20}^{2}} - \dfrac{81}{{b}^{2}} = 1, \\  \dfrac{{\left( {20}\sqrt{2}\right) }^{2}}{{20}^{2}} - \dfrac{{\left( t - {60}\right) }^{2}}{{b}^{2}} = 1, \end{array}\right. \text{ 解得 }\left\{  \begin{array}{l} b = {40}, \\  t = {20}, \end{array}\right.$
            
            所以双曲线 ${C}_{1}$ 的方程为 $\dfrac{{x}^{2}}{400} - \dfrac{{y}^{2}}{1600} = 1$ .
            
            \item 双曲线 ${C}_{1}$ 的共轭双曲线 ${C}_{2}$ 的方程为 $\dfrac{{y}^{2}}{1600} - \dfrac{{x}^{2}}{400} = 1$ .
            
            \item 双曲线 ${C}_{1}$ 与 ${C}_{2}$ 的离心率 ${e}_{1}$ 与 ${e}_{2}$ 满足关系式为 $\dfrac{1}{{e}_{1}^{2}} + \dfrac{1}{{e}_{2}^{2}} = 1$
            
            证明如下：
            
            $\dfrac{1}{{e}_{1}^{2}} + \dfrac{1}{{e}_{2}^{2}} = {\left( \dfrac{a}{c}\right) }^{2} + {\left( \dfrac{b}{c}\right) }^{2} = \dfrac{a^2+b^2}{c^2} = 1$
            
        \end{enumerate} 
    \end{jieda}
\end{enumerate}
\subsection*{双曲线的第二定义}
平面内到定点的距离和定直线的距离之比为定值$e$($e > 1$)的点的轨迹为双曲线. 这个定点就是双曲线的一个焦点, 这条定直线叫做双曲线的\textbf{准线},而常数 $e\left( {e > 1}\right)$ 是双曲线的离心率.
        
由双曲线的对称性可知,双曲线 $\dfrac{{x}^{2}}{{a}^{2}} - \dfrac{{y}^{2}}{{b}^{2}} = 1\left( {a > 0, b > 0}\right)$ 有两条准线: 与焦点 ${F}_{1}\left( {-c,0}\right)$ 的对应的准线方程是 $x =  - \dfrac{{a}^{2}}{c}$ ; 与焦点 ${F}_{2}\left( {c,0}\right)$ 对应的准线方程是 $x = \dfrac{{a}^{2}}{c}$ .

例题：
\begin{enumerate}
    \item 已知动点 $M\left( {x,y}\right)$ 到定点 ${F}_{1}\left( {-c,0}\right)$ 的距离和它到定直线 $l : x =  - \dfrac{{a}^{2}}{c}$ 的距离之比为常数 $\dfrac{c}{a}\left( {c > a > 0}\right)$ ,求动点 $M$ 的轨迹.

    \begin{jieda}
        设 $d$ 是动点 $M\left( {x,y}\right)$ 到直线 $l : x =  - \dfrac{{a}^{2}}{c}$ 的距离,则 $\dfrac{\left| M{F}_{1}\right| }{d} = \dfrac{c}{a}$ .
        
        即$\dfrac{\sqrt{{\left( x + c\right) }^{2} + {y}^{2}}}{\left| -\dfrac{{a}^{2}}{c} - x\right| } = \dfrac{c}{a}.$
        
        化简, 得$\left( {{a}^{2} - {c}^{2}}\right) {x}^{2} + {a}^{2}{y}^{2} = {a}^{2}\left( {{a}^{2} - {c}^{2}}\right) .$
        
        设 ${c}^{2} - {a}^{2} = {b}^{2}\left( {b > 0}\right)$ ,则有 $\dfrac{{x}^{2}}{{a}^{2}} - \dfrac{{y}^{2}}{{b}^{2}} = 1\left( {a > 0,b > 0}\right)$ . 这就是双曲线的标准方程,所以点 $M$ 的轨迹是双曲线.
    \end{jieda}
    \item 设点$M(m, n)$为双曲线$\dfrac{x^2}{a^2}-\dfrac{y^2}{b^2} = 1(a>0,b>0)$上的一点，$r_1, r_2$分别是$M$与双曲线左右焦点$F_1, F_2$的距离，试用双曲线离心率$e$和$M$的坐标表示$r_1, r_2$.
    \begin{jieda}
        根据双曲线的第二定义有$\dfrac{r_1}{\left|m - \left(-\dfrac{a^2}{c}\right)\right|} = \dfrac{r_2}{\left|m - \dfrac{a^2}{c}\right|} = e$，因此有$\left\{\begin{array}{l}
             r_1 = |a+em| \\
             r_2 = |a-em|
        \end{array}\right.$
        
        \ps 去绝对值：同负异正.
    \end{jieda}
    
    \item 过双曲线 ${x}^{2} - {y}^{2} = 4$ 的右焦点 $F$ 作倾斜角为 ${30}^{ \circ  }$ 的直线,交双曲线于 $A,B$ 两点,则弦长 $\left| {AB}\right|  =$ \blkda{8}.
    \begin{jieda}
        $e = \sqrt{2}$，$p = \dfrac{b^2}{c} = \dfrac{4}{2\sqrt{2}} = \sqrt{2}$，$|AB| = \left|\dfrac{2ep}{1-e^2cos^2\ang{30}}\right| = \left|\dfrac{4}{1-2\times \dfrac{3}{4}}\right| = 8$
    \end{jieda}
    
\end{enumerate}
\subsection*{双曲线的第三定义}
平面内的动点到两定点（设为$A, B$）的斜率乘积等于常数$m > 0$的点的轨迹是双曲线（刨除两定点）
\begin{dingautolist}{172}
    \item 当$m > k^2_{AB}$时，双曲线的实轴与$x$轴平行，$m = e^2 - 1$
    \item 当$k^2_{AB} > m > 0$时，双曲线的实轴与$y$轴平行，$m = \dfrac{1}{e^2-1}$
    \item 当$m = k^2_{AB}$时，双曲线退化为两相交直线.
\end{dingautolist}

特别地，不考虑$m = k^2_{AB}$的特殊情况下，根据双曲线离心率与渐近线斜率的关系可知，双曲线的第三定义可表述为：平面内的动点到两定点的斜率乘积等于常数$m > 0$的点的轨迹是实轴与坐标轴平行的双曲线（刨除两定点），且常数$m$为双曲线的渐近线的斜率平方.

证明：

设$A(x_1, y_1)$，$B(x_2, y_2)$，设$P(x, y)$，$k_{PA} = \dfrac{y-y_1}{x-x_1}$，$k_{PB} = \dfrac{y - y_2}{x-x_2}$

归一化（把对称中心移到原点）：

设$\left\{ \begin{array}{l}
     Y = y - \dfrac{y_1+y_2}{2}  \\
     X = x - \dfrac{x_1+x_2}{2}
\end{array}\right.$，$\left\{ \begin{array}{l}
     Y_1 = y_1 - \dfrac{y_1+y_2}{2} = \dfrac{y_1-y_2}{2}  \\
     X_1 = x_1 - \dfrac{x_1+x_2}{2} = \dfrac{x_1-x_2}{2}
\end{array}\right.$，$\left\{ \begin{array}{l}
     Y_2 = y_2 - \dfrac{y_1+y_2}{2} = -Y_1 \\
     X_2 = x_2 - \dfrac{x_1+x_2}{2} = -X_1
\end{array}\right.$

$m = k_{AP} \cdot k_{BP} = \dfrac{(y-y_1)(y-y_2)}{(x-x_1)(x-x_2)} = \dfrac{(Y-Y_1)(Y-Y_2)}{(X-X_1)(X-X_2)} = \dfrac{Y-Y_1^2}{X-X_1^2}$，即$Y^2- Y_1^2 = m(X^2-X_1^2)$

\begin{enumerate}
    \item $m = 0$，$Y = \pm Y_1$，即$y = y_1$或$y = y_2$.
    \item $m \neq 0$
    
    $Y^2- Y_1^2 = m(X^2-X_1^2) \Longleftrightarrow Y^2 - mX^2 = Y_1^2 - mX_1^2$
    \begin{dingautolist}{172}
        \item $m > \dfrac{Y_1^2}{X_1^2}$（即$m > k_{AB}^2$）

        $\dfrac{Y^2}{Y_1^2 - mX_1^2} - \dfrac{mX^2}{Y_1^2 - mX_1^2} = 1$

        $Y_1^2 - mX_1^2 < 0$，$\dfrac{Y^2}{Y_1^2 - mX_1^2} - \dfrac{mX^2}{Y_1^2 - mX_1^2} = 1$为实轴和$X$轴平行的双曲线方程.代换回$x, y$则是实轴与$x$轴平行的，以$\left(\dfrac{x_1+x_2}{2}, \dfrac{y_1+y_2}{2}\right)$为对称中心的双曲线.

        $a^2 = -\dfrac{Y_1^2 - mX_1^2}{m}$，$b^2 = -(Y_1^2 - mX_1^2)$，$c^2 = -(Y_1^2 - mX_1^2)\left(1+\dfrac{1}{m}\right)$，故$e^2 = m\left(1+\dfrac{1}{m}\right) = 1+m$
        \item $m = \dfrac{Y_1^2}{X_1^2}$（即$m = k_{AB}^2$）

        $Y = \pm \sqrt{m}X$为两相交直线.代换回$x, y$则是过$\left(\dfrac{x_1+x_2}{2}, \dfrac{y_1+y_2}{2}\right)$，斜率为$k_{AB}$和$-k_{AB}$的相交直线.

        \item $\dfrac{Y_1^2}{X_1^2} > m > 0$（即$0 < m < k_{AB}^2$）

        $\dfrac{Y^2}{Y_1^2 - mX_1^2} - \dfrac{mX^2}{Y_1^2 - mX_1^2} = 1$

        $Y_1^2 - mX_1^2 > 0$，$\dfrac{Y^2}{Y_1^2 - mX_1^2} - \dfrac{mX^2}{Y_1^2 - mX_1^2} = 1$为实轴和$Y$轴平行的双曲线方程.代换回$x, y$则实轴是与$y$轴平行的，以$\left(\dfrac{x_1+x_2}{2}, \dfrac{y_1+y_2}{2}\right)$为对称中心的双曲线.

        $b^2 = \dfrac{Y_1^2 - mX_1^2}{m}$，$a^2 = (Y_1^2 - mX_1^2)$，$c^2 = (Y_1^2 - mX_1^2)\left(1+\dfrac{1}{m}\right)$，故$e^2 = \left(1+\dfrac{1}{m}\right)$
        \item $0 > m > -1$

        $\dfrac{Y^2}{Y_1^2 - mX_1^2} - \dfrac{mX^2}{Y_1^2 - mX_1^2} = 1$为长轴和$X$轴平行的椭圆，代换回$x, y$则是长轴与$x$轴平行的，以$\left(\dfrac{x_1+x_2}{2}, \dfrac{y_1+y_2}{2}\right)$为对称中心的椭圆.

        $a^2 = -\dfrac{Y_1^2 - mX_1^2}{m}$，$b^2 = (Y_1^2 - mX_1^2)$，$c^2 = -(Y_1^2 - mX_1^2)\left(1+\dfrac{1}{m}\right)$，故$e^2 = m\left(1+\dfrac{1}{m}\right) = 1+m$
        
        \item $m = -1$

        $\dfrac{Y^2}{Y_1^2 - mX_1^2} - \dfrac{mX^2}{Y_1^2 - mX_1^2} = 1$为圆，代换回$x, y$则、以$\left(\dfrac{x_1+x_2}{2}, \dfrac{y_1+y_2}{2}\right)$为圆心的圆.
        
        \item $m < -1$

        $\dfrac{Y^2}{Y_1^2 - mX_1^2} - \dfrac{mX^2}{Y_1^2 - mX_1^2} = 1$为长轴和$Y$轴平行的椭圆，代换回$x, y$则是长轴与$y$轴平行的，以$\left(\dfrac{x_1+x_2}{2}, \dfrac{y_1+y_2}{2}\right)$为对称中心的椭圆.

        
        $b^2 = -\dfrac{Y_1^2 - mX_1^2}{m}$，$a^2 = (Y_1^2 - mX_1^2)$，$c^2 = (Y_1^2 - mX_1^2)\left(1+\dfrac{1}{m}\right)$，故$e^2 = \left(1+\dfrac{1}{m}\right)$
    \end{dingautolist}
\end{enumerate}

双曲线（椭圆）的第三定义说明双曲线（椭圆）的性质：在曲线上选取一组关于对称中心对称的点，那么双曲线上除两定点外的任意一点和他们连线的斜率乘积是定值.

例题：
\begin{enumerate}
    \item 过原点 $O$ 的直线 $l$ 与双曲线 $C : \dfrac{{x}^{2}}{{a}^{2}} - \dfrac{{y}^{2}}{{b}^{2}} = 1\left( {a > 0,b > 0}\right)$ 交于 $A,B$ 两点, $D$ 为 $C$ 的右顶点,直线 ${DA}$ 与直线 ${DB}$ 的斜率之积为3,则 $C$ 的渐近线方程为\blkda[2]{$y = \pm \sqrt{3}x$}
    \begin{jieda}
        根据双曲线的第三定义可知$e^2 - 1 = 3$

        根据双曲线渐近线斜率和离心率的关系可知$k^2 = e^2 - 1$，故$k = \pm \sqrt{3}$
    \end{jieda}
    \item 如图所示, $P$ 是双曲线 $\dfrac{{x}^{2}}{4} - {y}^{2} = 1$ 右支(在第一象限内)上的任意一点, ${A}_{1},{A}_{2}$ 分别是左右顶点, $O$ 是坐标原点,直线 $P{A}_{1},{PO},P{A}_{2}$ 的斜率分别为 ${k}_{1},{k}_{2},{k}_{3}$ ,则斜率之积 ${k}_{1}{k}_{2}{k}_{3}$ 的取值范围是\blkda{$\left(0, \dfrac{1}{8}\right)$}.
    \begin{center}
        \includegraphics[width=0.2\linewidth]{bodyxb//src/shuangquxiandisandingyi.png}
    \end{center}
    \begin{jieda}
        $k_1k_2 = \dfrac{1}{4}$
    \end{jieda}
    \item 若双曲线 ${x}^{2} - {y}^{2} = {a}^{2}\left( {a > 0}\right)$ 的左、右顶点分别为 $A\text{ 、 }B$ ,点 $P$ 是第一象限内双曲线上的点. 若直线 ${PA}\text{ 、 }{PB}$ 的倾斜角分别为 $\alpha ,\beta$ ,那么 $\alpha + \beta = $\blkda{$\dfrac{\pi}{2}$}
    \begin{jieda}
        等轴双曲线的渐近线斜率$k = \pm 1$，故根据双曲线第三定义可知$PA, PB$斜率乘积为$e^2-1 = k^2 = 1$，即$PA, PB$斜率互为倒数，根据正切和余切的商数关系和诱导公式可知两直线的倾斜角互余.
    \end{jieda}
    \item 如图所示,已知 $M$ , $N$ 为双曲线 $\dfrac{{x}^{2}}{{a}^{2}} - \dfrac{{y}^{2}}{{b}^{2}} = 1$ ( $a > 0$ , $b > 0$ )上关于原点对称的两点,点 $M$ 与点 $Q$ 关于 $x$ 轴对称, $Q$为$ME$中点,直线 ${NE}$ 交双曲线右支于点 $P$ ,若 $\angle {NMP} = \dfrac{\pi }{2}$ ,则 $e =$\blkda{$\sqrt{2}$}
    \begin{center}
        \includegraphics[width=0.2\linewidth]{bodyxb//src/disandingyi3.png}
    \end{center}
    \begin{jieda}
        设$M(x_0, y_0)$，则$N(-x_0, -y_0)$，$Q(x_0, -y_0)$，$E(x_0, -3y_0)$，故$k_{NP} = -\dfrac{y_0}{x_0}, k_{MN} = \dfrac{y_0}{x_0}, k_{MP} = -\dfrac{x_0}{y_0}$，根据第三定义有$1 = k_{NP} \cdot k_{MP} = e^2 - 1$，故$e =\sqrt{2}$
    \end{jieda}
\end{enumerate}
\subsection*{点和双曲线的位置关系}
$\dfrac{x^2}{a^2} - \dfrac{y^2}{b^2} = 1$双曲线及其渐近线$y = \pm \dfrac{b}{a}x$将空间分成五个区域，对于点$(x_0, y_0)$设$t = \dfrac{x_0^2}{a^2} - \dfrac{y_0^2}{b^2}$，则有$
\begin{cases}
t < 0 & \text{点在渐近线之间无双曲线的区域内}, \\
t = 0 & \text{点在渐近线上}, \\
0 < t < 1 & \text{点在渐近线与双曲线之间的区域内}, \\
t = 1 & \text{点在双曲线上}, \\
t > 1 & \text{点在双曲线内}.
\end{cases}$
\subsection*{直线和双曲线的位置关系}
直线与双曲线的位置关系可以通过代数方法进行判断.
\begin{enumerate}
    \item 联立:把直线与双曲线的方程联立,构成二元二次方程组
    \item 消元:消去未知数 $x$ (或 $y$ ),得到关于 $y$ (或 $x$ )的一元二次方程（\textbf{二次项系数为0对应着直线与双曲线的渐近线平行}）
    \item 计算该一元二次方程的判别式 $\Delta$ ,并得出结论:
    \begin{dingautolist}{172}
        \item $\Delta  > 0 \Longleftrightarrow$ 有两个公共点
        \item $\Delta  = 0 \Longleftrightarrow$ 有且只有一个公共点
        \item $\Delta  < 0 \Longleftrightarrow$ 没有公共点
    \end{dingautolist}
\end{enumerate}
直线$l$的斜率$k$和双曲线$C: \dfrac{{x}^{2}}{{a}^{2}} - \dfrac{{y}^{2}}{{b}^{2}} = 1\left( {a > 0, b > 0}\right)$的渐近线的斜率$\pm k'$大小关系与$l$与$C$交点个数之间的关系
\begin{enumerate}
    \item $|k| < |k'|$时与双曲线的两支各有一个交点，共两个交点
    \item $|k| = |k'|$时有一个交点（直线与渐近线重合时没有交点）
    \item $|k| > |k'|$时与双曲线的一支有两个交点或相切或与双曲线没有交点
\end{enumerate}

切线结论:
\begin{enumerate}
    \item 双曲线$C: \dfrac{{x}^{2}}{{a}^{2}} - \dfrac{{y}^{2}}{{b}^{2}} = 1\left( {a > 0, b > 0}\right)$在$(x_0, y_0)$处的切线为$\dfrac{{x}x_0}{{a}^{2}} - \dfrac{{y}y_0}{{b}^{2}} = 1$
    \item 过点$M(x_0, y_0)$（在双曲线外且不在渐近线上），做双曲线$\dfrac{x^2}{a^2}-\dfrac{y^2}{b^2} = 1({a > 0, b > 0})$的切线$l_1, l_2$，切点为$A, B$，则直线$AB$的方程为$\dfrac{x_0x}{a^2}-\dfrac{y_0y}{b^2} = 1$
\end{enumerate}

例题：
\begin{enumerate}
    \item 设 $k$ 为实数,已知双曲线 $C$ 的方程为 $\dfrac{{x}^{2}}{4} - {y}^{2} = 1$ ,直线 $l$ 的方程是 $y = {kx} + 1$ . 当实数 $k$ 为何值时,直线 $l$ 与双曲线 $C$ 满足:
    \begin{dingautolist}{172}
        \item 有两个公共点
        \item 仅有一个公共点
        \item 没有公共点
    \end{dingautolist}
    \begin{jieda}
        联立方程组 $\left\{  \begin{array}{l} \dfrac{{x}^{2}}{4} - {y}^{2} = 1, \\  y = {kx} + 1. \end{array}\right.$ 消去 $y$，整理得$\left( {1 - 4{k}^{2}}\right) {x}^{2} - {8kx} - 8 = 0.$
        \begin{dingautolist}{172}
            \item 当直线 $l$ 与双曲线 $C$ 有两个公共点时,有 $\left\{  \begin{array}{l} 1 - 4{k}^{2} \neq  0, \\  \Delta  = {64}{k}^{2} + {32}\left( {1 - 4{k}^{2}}\right)  > 0. \end{array}\right.$
            
            解得 $k \in  \left( {-\dfrac{\sqrt{2}}{2}, - \dfrac{1}{2}}\right)  \cup  \left( {-\dfrac{1}{2},\dfrac{1}{2}}\right)  \cup  \left( {\dfrac{1}{2},\dfrac{\sqrt{2}}{2}}\right)$ .
            \item 当直线 $l$ 与双曲线 $C$ 仅有一个公共点时,有 $1 - 4{k}^{2} = 0$ 或 $\left\{  \begin{array}{l} 1 - 4{k}^{2} \neq  0, \\  \Delta  = {64}{k}^{2} + {32}\left( {1 - 4{k}^{2}}\right)  = 0. \end{array}\right.$ 解得 $k =  \pm  \dfrac{1}{2}$ 或 $\pm  \dfrac{\sqrt{2}}{2}$ .
            \item 当直线 $l$ 与双曲线 $C$ 没有公共点时,有 $\left\{  \begin{array}{l} 1 - 4{k}^{2} \neq  0, \\  \Delta  = {64}{k}^{2} + {32}\left( {1 - 4{k}^{2}}\right)  < 0. \end{array}\right.$ 解得 $k \in  \left( {-\infty , - \dfrac{\sqrt{2}}{2}}\right)  \cup  \left( {\dfrac{\sqrt{2}}{2}, + \infty }\right)$ .
        \end{dingautolist}
    \end{jieda}
    \item 已知直线 $l$ 的方程为 $y = {kx} - 1$ ,双曲线 $C$ 的方程为 ${x}^{2} - {y}^{2} = 1$ . 若直线 $l$ 与双曲线 $C$ 的右支交于不同的两点,则实数 $k$ 的取值范围是\blkda{$\left( {1,\sqrt{2}}\right)$}
    \begin{jieda}
        \begin{enumerate}
            \item 方法1 

            联立 $\left\{  \begin{array}{l} y = {kx} - 1, \\  {x}^{2} - {y}^{2} = 1, \end{array}\right.$，整理得 $\left( {1 - {k}^{2}}\right) {x}^{2} + {2kx} - 2 = 0$，因为直线 $y = {kx} - 1$ 与双曲线 ${x}^{2} - {y}^{2} = 1$ 的右支交于不同的两点，所以 $\left\{  \begin{array}{l} 1 - {k}^{2} \neq  0, \\  \Delta  = 4{k}^{2} + 8\left( {1 - {k}^{2}}\right)  > 0, \\  \dfrac{-{2k}}{1 - {k}^{2}} > 0, \\  \dfrac{-{2}}{1 - {k}^{2}} > 0, \end{array}\right.$
            
            解得 $1 < k < \sqrt{2}$，所以实数 $k$ 的取值范围为 $\left( {1,\sqrt{2}}\right)$ .

            \item 方法2 

             $y = {kx} - 1$恒过定点$(0， -1)$，易知当$k = 1$时和双曲线的渐近线平行，而当直线斜率增大的过程中会先与双曲线右支交于两点，而后达到相切的状态，因此只需再求得相切时斜率的右边界即可.

             联立 $\left\{  \begin{array}{l} y = {kx} - 1, \\  {x}^{2} - {y}^{2} = 1, \end{array}\right.$，整理得 $\left( {1 - {k}^{2}}\right) {x}^{2} + {2kx} - 2 = 0$，令$\Delta = 0$，即$4{k}^{2} + 8\left( {1 - {k}^{2}}\right) = 0$，得到$k = \pm \sqrt{2}$，结合上述分析可知右边界为$\sqrt{2}$

             因此$k$ 的取值范围为 $\left( {1,\sqrt{2}}\right)$ .
        \end{enumerate}
        
    \end{jieda}
    \item 如图,已知 $M$ 为双曲线 $E : \dfrac{{x}^{2}}{{a}^{2}} - \dfrac{{y}^{2}}{{b}^{2}} = 1\left( {a > 0,b > 0}\right)$ 上一动点,过 $M$ 作双曲线 $E$ 的切线交 $x$ 轴于点 $A$ ,过点 $A$ 作 ${AD}\bot {OM}$ 于点 $D,\left| {OD}\right|  \cdot  \left| {OM}\right|  = 2{b}^{2}$ ,则双曲线 $E$ 的离心率为\blkda{$\dfrac{\sqrt{6}}{2}$}
    
    \begin{center}
    \includegraphics[width=0.3\textwidth]{bodyxb/src/bo_d3mggojef24c73cuv41g_0_235_629_430_302_0.jpg}
    \end{center}

    \begin{jieda}
        如图所示，根据$\triangle OAD \sim \triangle OMB$可知$\dfrac{|OD|}{|OB|} = \dfrac{|OA|}{|OM|}$，故$|OA| \cdot |OB| = 2b^2$
        \begin{center}
        \includegraphics[width=0.3\textwidth]{bodyxb/src/bo_d3mggojef24c73cuv41g_3_236_262_430_302_0.jpg}
        \end{center}

        设$M(x_0, y_0)$，则$AM: \dfrac{x_0x}{a^2}-\dfrac{y_0y}{b^2} = 1$，故$|OA| = \dfrac{a^2}{x_0}$，因此$|OA| \cdot |OB| = a^2 = 2b^2$，故$e = \dfrac{\sqrt{6}}{2}$
    \end{jieda}
    
    \item 已知双曲线 ${C}$ 的方程为 ${x}^{2} - \dfrac{{y}^{2}}{2} = 1$ ,试判断过点 $P\left( {1,1}\right)$ 能否作一条直线 $l$ 与双曲线交于 $A\text{ 、 }B$ 两点,且 $P$ 为 ${AB}$ 的中点,并说明理由.
    \begin{jieda}
        设双曲线上两点$A(x_1, y_1), B(x_2, y_2)$，$P$ 为 ${AB}$ 的中点，则有$\left\{\begin{array}{l}
             {x_1}^{2} - \dfrac{{y_1}^{2}}{2} = 1  \\
             {x_2}^{2} - \dfrac{{y_2}^{2}}{2} = 1
        \end{array}\right.$，两式作差有$(x_1+x_2)(x_1-x_2) - \dfrac{(y_1+y_2)(y_1-y_2)}{2} = 0$，根据$P$为中点有$x_1+x_2 = 2, y_1+y_2 = 2$，故$2(x_1-x_2) - (y_1-y_2) = 0$，即$k_{AB} = 2$，故$AB: y = 2x-1$

        联立直线和双曲线方程有$\left\{\begin{array}{l}
             y = 2x-1  \\
             {x}^{2} - \dfrac{{y}^{2}}{2} = 1
        \end{array}\right.$，消去$y$可得$2x^2-(2x-1)^2 = 2$，即$-2x^2+4x-3=0$，方程无解，故不存在.
    \end{jieda}

    \item 已知 $F\left( {2,0}\right)$ 为双曲线 $C : \dfrac{{x}^{2}}{{a}^{2}} - \dfrac{{y}^{2}}{{b}^{2}} = 1\left( {a > 0,b > 0}\right)$ 的右焦点,且点 $F$ 到双曲线 $C$ 的一条渐近线的距离为 $\sqrt{3}$ .
    \begin{enumerate}
        \item 求双曲线 $C$ 的方程
        \item 设过点 $F$ 的直线 $l$ 与双曲线 $C$ 相交于点 $P$ 、 $Q$ ,线段 ${PQ}$ 的垂直平分线与 $x$ 轴交于点 $M$ ,若 $\left| {MF}\right|  = 6$ ,求直线 $l$ 的方程.
    \end{enumerate}
    \begin{jieda}
        \begin{enumerate}
            \item 因为 ${bx} - {ay} = 0$ 是双曲线 $C$ 的一条渐近线方程,则 $\dfrac{\left| 2b\right| }{\sqrt{{a}^{2} + {b}^{2}}} = \sqrt{3}$ ,所以 $b = \sqrt{3}$ , 又因为 $c = 2$ ,所以 ${a}^{2} = {c}^{2} - {b}^{2} = 1$ ,即 $a = 1$ .
            
            所以双曲线 $C$ 的方程为 ${x}^{2} - \dfrac{{y}^{2}}{3} = 1$ .
            \item 
            设$P\left( {{x}_{1},{y}_{1}}\right)$ 、 $Q\left( {{x}_{2},{y}_{2}}\right)$，其中点为$N(x_0, y_0)$则$\left\{\begin{array}{l}
                 x_1^{2} - \dfrac{y_1^{2}}{3} = 1 \\
                 x_2^{2} - \dfrac{y_2^{2}}{3} = 1
            \end{array}\right.$
            两式做差可得$(x_1+x_2)(x_1-x_2) - \dfrac{(y_1+y_2)(y_1-y_2)}{3} = 0$，易知$k_{PQ}$存在，故$k_{PQ} = \dfrac{y_1-y_2}{x_1-x_2} = \dfrac{3(x_1+x_2)}{y_1+y_2} = \dfrac{3x_0}{y_0}$，故$PQ$为$y = \dfrac{3x_0}{y_0}(x-2)$，代入$N(x_0, y_0)$有$y_0^2= 3x_0^2-6x_0$
            
            $PQ$的垂直平分线为$y = -\dfrac{1}{k_{PQ}}(x-x_0)+y_0 = -\dfrac{y_0}{3x_0}(x-x_0)+y_0 = -\dfrac{y_0}{3x_0}x+\dfrac{4y_0}{3}$
            \begin{dingautolist}{172}
                \item $PQ$的垂直平分线过$(8, 0)$，则有$-\dfrac{8y_0}{3x_0}+\dfrac{4y_0}{3} = 0$，故$x_0 = 2$，根据$y_0^2= 3x_0^2-6x_0$可知$y_0 = 0$(舍)
                \item $PQ$的垂直平分线过$(-4, 0)$，则有$\dfrac{4y_0}{3x_0}+\dfrac{4y_0}{3} = 0$，故$x_0 = -1$，根据$y_0^2= 3x_0^2-6x_0$可知$y_0 = \pm 3$，故直线 $l$ 的方程为 $x + y - 2 = 0$ 或 $x - y - 2 = 0$ .
            \end{dingautolist}
        \end{enumerate}
    \end{jieda}

    \item 已知双曲线 $\Gamma  : {x}^{2} - {y}^{2} = 1$ ,直线 $l$ 过点 $M\left( {m,0}\right) ,m \in  \mathbf{R}$ ,且与 $\Gamma$ 的右支交于 $P\text{、}Q$ 两点，与 $\Gamma$ 的两条渐近线分别交于 $A$ 、 $B$ 两点,其中 $A$ 、 $P$ 在第一象限, $B$ 、 $Q$ 在第四象限.
    \begin{enumerate}
        \item 求 $\Gamma$ 的两条渐近线的夹角
        \item 若 $M$ 为 $\Gamma$ 的右焦点,求 $\bigtriangleup {OAB}$ 的面积的最小值
        \item 若 $\left| {AP}\right|  = \left| {PQ}\right|  = \left| {QB}\right|$ ,求 $m$ 的取值范围.
    \end{enumerate}
    \begin{jieda}
        \begin{enumerate}
            \item 两条渐近线为$y = \pm x$，故夹角为$\dfrac{\pi}{2}$
            \item 当直线$l$斜率为0时不符合要求，故设直线$l: x = ty+m$，由$l$与 $\Gamma$ 的右支交于 $P\text{、}Q$两点，且$P$在第一象限，$Q$在第四象限，易知$t \in (-1, 1), m > 1$（第三问相同约束）
            
            联立两条渐近线和直线$l$：$\left\{\begin{array}{l}
                x^2-y^2 = 0 \\
                x = ty+m
            \end{array} \right.$，$(t^2-1)y^2+2tmy+m^2 = 0$，必有$\Delta > 0$，故$\left\{\begin{array}{l}
                y_A + y_B = \dfrac{-2tm}{t^2-1} \\
                y_Ay_B = \dfrac{m^2}{t^2-1}
            \end{array}\right.$

            $S_{\triangle OAB} = S_{\triangle MOA} + S_{\triangle MOB} = \dfrac{m}{2} \cdot |y_A - y_B| = \dfrac{|m|}{2} \cdot \dfrac{\sqrt{4t^2m^2 - 4m^2(t^2-1)}}{t^2-1} = \dfrac{m^2}{1-t^2}$

            因为$l$过右焦点，所以$m = \sqrt{2}$，因此$S_{\triangle OAB} = \dfrac{2}{1-t^2} \leq 2$（$t = 0$）

            \item 联立两条渐近线和直线$l$：$\left\{\begin{array}{l}
                x^2-y^2 = 1 \\
                x = ty+m
            \end{array} \right.$，$(t^2-1)y^2+2tmy+m^2 - 1 = 0$，必有$\Delta > 0$，故$\left\{\begin{array}{l}
                y_P + y_Q = \dfrac{-2tm}{t^2-1} \\
                y_Py_Q = \dfrac{m^2-1}{t^2-1}
            \end{array}\right.$
            
            因为$y_P + y_Q = y_A + y_B$，所以$|AB|, |PQ|$共中点，因此$\left| {AP}\right|  = \left| {PQ}\right|  = \left| {QB}\right|$即$|AB| = 3|PQ|$，即$|y_A-y_B| = 3|y_P - y_Q|$，即$\dfrac{\sqrt{4t^2m^2 - 4m^2(t^2-1)}}{t^2-1} = 3\dfrac{\sqrt{4t^2m^2 - 4(m^2-1)(t^2-1)}}{t^2-1}$，即$\sqrt{4m^2}  = 3\sqrt{4(m^2+t^2-1)}$，即$m^2 = \dfrac{9(1-t^2)}{8} \in \left(0, \dfrac{9}{8}\right]$，又因为$m > 1$，所以$m \in \left(1, \dfrac{3\sqrt{2}}{4}\right]$
        \end{enumerate}
    \end{jieda}
    
    \item 已知双曲线 $\Gamma  : \dfrac{{x}^{2}}{2} - \dfrac{{y}^{2}}{2} = 1$ ,过点 $M\left( {m,0}\right)$ 作不垂直于 $x$ 轴的直线 $l$ 交双曲线 $\Gamma$ 于 $A$ 、 $B$ 两点.
    \begin{enumerate}
        \item 求双曲线离心率
        \item 若点 $A\left( {\sqrt{3},1}\right)$ ,点 $B$ 在双曲线的右支上,且 $B$ 是 ${AM}$ 的中点,求直线 $l$ 的斜率
        \item 若 $m > 0,{F}_{1},{F}_{2}$ 分别是双曲线 $\Gamma$ 的左右焦点, ${A}^{\prime }$ 是 $A$ 关于 $y$ 轴的对称点,若存在直线 $l$使得 $\vv{{F}_{1}{A}^{\prime }} \cdot  \vv{{F}_{2}B} = 0$ ,求 $m$ 的取值范围.
    \end{enumerate}
    \begin{jieda}
        \begin{enumerate}
            \item $c = 2, a = \sqrt{2}, e = \dfrac{c}{a} = \sqrt{2}$
            \item $B\left(\dfrac{\sqrt{3}+m}{2}, \dfrac{1}{2}\right)$在双曲线上，$\left(\dfrac{\sqrt{3}+m}{2}\right)^2 - \dfrac{1}{4} = 2$，故$(\sqrt{3}+m)^2 =9$又因为$B$在右支上，所以$\dfrac{\sqrt{3}+m}{2} > 0$，因此$m = 3-\sqrt{3}$，即$B\left(\dfrac{3}{2}, \dfrac{1}{2}\right)$，因此$k_l = \dfrac{1-\dfrac{1}{2}}{\sqrt{3}-\dfrac{3}{2}} = \dfrac{1}{2\sqrt{3}-3} = \dfrac{2\sqrt{3}+3}{3}$
            \item $F_1(-2, 0)$，$F_2(2, 0)$
            
            当$l$为$y = 0$时不满足要求，故设直线$l:  x = ty + m$，设$A(x_1, y_1), B(x_2, y_2)$，因此$A'(-x_1, y_1)$

            $\vv{F_1A'} = (2-x_1, y_1)$，$\vv{F_2A} = (x_2-2, y_2)$，$\vv{F_1A'} \cdot \vv{F_2A} =(2-x_1)(x_2-2)+y_1y_2 = 0$，即$(2-m-ty_1)(ty_2+m-2)+y_1y_2 = 0$，即$(1-t^2)y_1y_2+(2-m)t(y_1+y_2)-(m-2)^2 = 0$

            联立直线与双曲线$\left\{\begin{array}{l}
                 x = ty+m  \\
                 x^2-y^2 = 2
            \end{array}\right.$，消去$x$得$(ty+m)^2-y^2 = 2$，即$(t^2-1)y^2+2tmy+(m^2-2) = 0$，$t \neq \pm 1$，$\Delta = 4t^2m^2 - 4(m^2-2)(t^2-1) = 4m^2+8t^2-8 > 0$

            根据韦达定理有$\left\{\begin{array}{l}
                 y_1+y_2 = \dfrac{2tm}{1-t^2} \\
                 y_1y_2 = \dfrac{2-m^2}{1-t^2}
            \end{array}\right.$，代入$(1-t^2)y_1y_2+(2-m)t(y_1+y_2) -(m-2)^2 = 0$得$(1-t^2)\cdot\dfrac{2-m^2}{1-t^2}+(2-m)t\cdot\dfrac{2tm}{1-t^2}-(m-2)^2 = 0$，即$\dfrac{t^2-m^2+2m-1}{1-t^2} = 0$，故$t^2 = (m-1)^2$

            代入$4m^2+8t^2-8 > 0$可得$m^2+2(m-1)^2-2>0$，解得$m < 0$或$m > \dfrac{4}{3}$，又因为$m > 0, t^2 \neq 1$，所以$m \in \left(\dfrac{4}{3}, 2\right) \cup (2, +\infty)$
        \end{enumerate}
    \end{jieda}

    \item 定义 $\Gamma  : {x}^{2} - {y}^{2} = 1,{F}_{1}\text{、}{F}_{2}$ 为 $\Gamma$ 的左右焦点 
    \begin{enumerate}
        \item 求点$(2,0)$到 $\Gamma$ 渐近线的距离
        \item ${P}$ 为 $\Gamma$ 上一点, $\vv{P{F}_{1}} \cdot  \vv{P{F}_{2}} = 1$ ,求 $\bigtriangleup P{F}_{1}{F}_{2}$ 的面积
        \item 设 $\Omega  : {x}^{2} - {y}^{2} = 1$ ,其中 $\left\{  \begin{array}{l} x < 0 \\  y \leq   - 1 \end{array}\right.$ 或 $\left\{  \begin{array}{l} x > 0 \\  y \geq   - 1 \end{array}\right.$ ,过点 ${F}_{2}$ 的直线 $l$ 交 $\Omega$ 于 $P\text{、}Q$ 两点(分别位于一、四象限),过点 ${F}_{2}$ 直线 $m$ 交$\Omega$ 于 $M\text{、}N$ 两点 (分别位于三、四象限),是否存在正数 $\lambda$ ,对于任意的 $l$ ,都存在唯一的 $m$ ,使 $\left| {MN}\right|  = \lambda \left| {PQ}\right|$ 成立,若存在,求出所有的 $\lambda$ ,若不存在,请说明理由.
    \end{enumerate}
    \begin{jieda}
        \begin{enumerate}
            \item 渐近线为$y = \pm x$，$(2, 0)$到渐近线的距离为$d = \sqrt{2}$
            \item $F_1(-\sqrt{2}, 0), F_2(\sqrt{2}, 0)$
            \begin{dingautolist}{172}
                \item 设$P(x_0, y_0)$，$\vv{PF_1} \cdot \vv{PF_2} = |\vv{OP}|^2 - 2 = 1$，故$\left\{\begin{array}{l}
                     x_0^2 + y_0^2 = 3 \\
                     x_0^2 - y_0^2 = 1
                \end{array}\right.$，故$y_0 = \pm 1$
                \item 设$P\left(\dfrac{1}{cos\theta}, tan\theta\right), \theta \in \left[0, 2\pi\right)$且$\theta \neq \dfrac{\pi}{2}, \dfrac{3\pi}{2}$，则$\vv{PF_1} = \left(-\sqrt{2} - \dfrac{1}{cos\theta}, -tan\theta\right)$，$\vv{PF_2} = \left(\sqrt{2} - \dfrac{1}{cos\theta}, -tan\theta\right)$，$1 = \vv{PF_1} \cdot \vv{PF_2} = \dfrac{1}{cos^2\theta} - 2 + tan^2\theta$，即$3cos^2\theta = 1 + sin^2\theta$，即$cos^2\theta = \dfrac{1}{2}$，故$y_P = \pm 1$
            \end{dingautolist}
            故$S_{\triangle{PF_1F_2}} = \sqrt{2}$
            
            \item 由题可知$x_M < 0 < x_N$，$x_Q < \sqrt{2} < x_P$，$k_l \in (1, +\infty)$，或$k_l$不存在，$k_m \in \left[\dfrac{\sqrt{2}}{4}, 1\right)$
            
            $|F_2M| = \sqrt{(x_M-\sqrt{2})^2+y_M^2} = -\sqrt{2}x_M+1$，同理$|F_2N| = \sqrt{2}x_N - 1$，故$|MN| = 2-\sqrt{2}(x_M+x_N)$

            $|F_2P| = \sqrt{2}x_P - 1$，$|F_2Q| = \sqrt{2}x_Q - 1$，故$|PQ| = \sqrt{2}(x_P + x_Q) - 2$

            设直线$PQ: y = k_l(x-\sqrt{2})$，与双曲线联立得$(1-k_l^2)x^2+2\sqrt{2}k_l^2x-(2k_l^2+1)=0$

            $\Delta>0$, $x_P+x_Q=-\dfrac{2\sqrt{2}k_l^2}{1-k_l^2}$，故$|PQ|= \dfrac{-4k_l^2}{1-k_l^2}-2 \in (2,+\infty)$, 当$k_l$不存在时，$|PQ|=2$，故$|PQ|\in [2,+\infty)$

            同理可得$|MN|=2+\dfrac{4k_m^2}{1-k_m^2}$，故$|MN|$随$k_m$增大严增，$|MN|\in \left[\dfrac{18}{7},+\infty\right)$

            因此当$\lambda \geq \dfrac{9}{7}$满足要求.
        \end{enumerate}
    \end{jieda}
    
    \item 已知双曲线 $C : \dfrac{{x}^{2}}{{a}^{2}} - \dfrac{{y}^{2}}{{b}^{2}} = 1\left( {a > 0,b > 0}\right)$ 的左、右顶点分别为 $P,Q$ ,右焦点 $F$ 到渐近线的距离为 1,且离心率为 $\dfrac{2\sqrt{3}}{3}$.
    \begin{enumerate}
        \item 求双曲线 $C$ 的标准方程
        \item 过点 $F$ 的直线 $l$ (直线 $l$ 的斜率不为 0 $)$ 与双曲线 $C$ 交于 $A$ ， $B$ 两点，若 $M$ ， $N$ 分别为直线 ${AP},{BQ}$ 与 $y$ 轴的交点,记 $\bigtriangleup {OPM},\bigtriangleup {OQN}$ 的面积分别记为 ${S}_{\bigtriangleup {OPM}}$ , ${S}_{\bigtriangleup {QQN}}$ ,求 $\dfrac{{S}_{\bigtriangleup {OPM}}}{{S}_{\bigtriangleup {OQN}}}$ 的值.
    \end{enumerate}
    \begin{jieda}
        \begin{enumerate}
            \item $\dfrac{x^2}{3} - y^2 = 1$
            \item $\dfrac{S_{\bigtriangleup {OPM}}}{S_{\bigtriangleup {OQN}}} = \dfrac{|OP| \cdot |OM|}{|OQ| \cdot |ON|} = \left|\dfrac{y_M}{y_N}\right|$
            \begin{center}
                \includegraphics[width=0.3\textwidth]{bodyxb/src/0_1186_533_385_309_0.jpg}
            \end{center}
            $F\left( {2,0}\right)$ ,设直线 $l : x = {my} + 2,A\left( {{x}_{1},{y}_{1}}\right) ,B\left( {{x}_{2},{y}_{2}}\right)$
            联立 $\left\{  \begin{array}{l} \dfrac{{x}^{2}}{3} - {y}^{2} = 1 \\  x = {my} + 2 \end{array}\right.$ ,得 $\left( {{m}^{2} - 3}\right) {y}^{2} + {4my} + 1 = 0,{m}^{2} - 3 \neq  0$,
            $\left\{\begin{array}{l}
                {y}_{1} + {y}_{2} = \dfrac{-{4m}}{{m}^{2} - 3} \\
                {y}_{1}{y}_{2} = \dfrac{1}{{m}^{2} - 3}
            \end{array}\right. $

            直线 ${AP}$ 的方程为 $y = \dfrac{{y}_{1}}{{x}_{1} + \sqrt{3}}\left( {x + \sqrt{3}}\right)$ ,当 $x = 0$ 时, $y = \dfrac{\sqrt{3}{y}_{1}}{{x}_{1} + \sqrt{3}} = \dfrac{\sqrt{3}{y}_{1}}{m{y}_1 + 2 + \sqrt{3}}$

            直线 ${BQ}$ 的方程为 $y = \dfrac{{y}_{2}}{{x}_{2} - \sqrt{3}}\left( {x - \sqrt{3}}\right)$ ,当 $x = 0$ 时, $y = \dfrac{-\sqrt{3}{y}_{2}}{{x}_{2} - \sqrt{3}} = \dfrac{\sqrt{3}{y}_2}{m{y}_2 + 2 + \sqrt{3}}$

            即 $M\left(0, \dfrac{\sqrt{3}{y}_{1}}{m{y}_1 + 2 + \sqrt{3}}\right) ,N\left(0, \dfrac{\sqrt{3}{y}_2}{m{y}_2 + 2 + \sqrt{3}}\right)$

            $\dfrac{{S}_{\bigtriangleup {OPM}}}{{S}_{\bigtriangleup {OQN}}} = \left| \dfrac{{y}_{M}}{{y}_{N}}\right| = \left| \dfrac{{y}_{l}\left( {m{y}_{2} + 2 - \sqrt{3}}\right) }{{y}_{2}\left( {m{y}_1 + 2 + \sqrt{3}}\right) }\right| = \left| \dfrac{m{y}_{1}{y}_{2} + \left( {2 - \sqrt{3}}\right) {y}_{1}}{m{y}_{1}{y}_{2} + \left( {2 + \sqrt{3}}\right) {y}_{2}}\right|$
            
            因为$m{y}_{1}{y}_{2} = -\dfrac{1}{4}(y_1+y_2)$，所以$\dfrac{{S}_{\bigtriangleup {OPM}}}{{S}_{\bigtriangleup {OQN}}} = \left| \dfrac{-\dfrac{1}{4}(y_1+y_2) + \left( {2 - \sqrt{3}}\right) {y}_{1}}{-\dfrac{1}{4}(y_1+y_2) + \left( {2 + \sqrt{3}}\right) {y}_{2}}\right| = \left|\dfrac{\left(\dfrac{7}{4}-\sqrt{3}\right)y_1-\dfrac{1}{4}y_2}{-\dfrac{1}{4}y_1+\left(\dfrac{7}{4}+\sqrt{3}\right)y_2}\right|= 7 - 4\sqrt{3}$.
        \end{enumerate}
    \end{jieda}

    \item 设 $a > 0,b > 0,m\text{、}t \in  \mathbf{R}$ ,双曲线 $\Gamma  : \dfrac{{x}^{2}}{{a}^{2}} - \dfrac{{y}^{2}}{{b}^{2}} = 1$ 的一条渐近线方程是 $y = 2\sqrt{2}x$ ,点 $P$为 $\Gamma$ 右支上的一点,直线 $l$ 的方程是 $x - {my} - t = 0,O$ 是坐标原点.
    \begin{enumerate}
        \item 若点 $P$ 的坐标为(3,8),求双曲线 $\Gamma$ 的方程
        \item 若直线 $l$ 经过点 $O$，且与 $\Gamma$ 交于 $A$ 、 $B$ 两点,直线 ${PA}$ 、 ${PB}$ 的斜率分别为 ${k}_{1}$ 、 ${k}_{2}$，求 ${k}_{1} \cdot  {k}_{2}$ 的值
        \item 设点 ${F}_{1}$ 是 $\Gamma$ 的左焦点,点 ${A}_{1}$ 、 ${A}_{2}$ 是 $\Gamma$ 的左、右两个顶点,直线 ${A}_{1}P$ 与直线 $x = 1$ 交于点 $M$ ,直线 $l$ 经过点 $P$ 与 $\Gamma$ 的右支交于另外一点 $Q$ ,若 $a = 3$ ,且直线 ${MQ}$ 恒过点 ${A}_{2}$ ,求$\bigtriangleup {F}_{1}{PQ}$ 周长的取值范围.
    \end{enumerate}
    \begin{jieda}
        \begin{enumerate}
            \item $\dfrac{x^2}{a^2} - \dfrac{y^2}{8a^2} = 1$，代入$P(3, 8)$，$\dfrac{9}{a^2} - \dfrac{64}{8a^2} = 1$，解得$a = 1$，故$x^2 - \dfrac{y^2}{8} = 1$.
            \item \ps 背景为第三定义$k_1k_2 = e^2 - 1 = k^2 = 8$

            设$A(x_0, y_0), B(-x_0, -y_0), P(x_p, y_p)$，故$k_1 = \dfrac{y_p - y_0}{x_p - x_0}, k_2 = \dfrac{y_p + y_0}{x_p + x_0}$，$k_1k_2 = \dfrac{y_p^2-y_0^2}{x_p^2-x_0^2}$

            $\dfrac{x^2}{a^2} - \dfrac{y^2}{8a^2} = 1$即$x^2 = a^2 + \dfrac{y^2}{8}$

            $k_1k_2 = \dfrac{y_p^2-y_0^2}{x_p^2-x_0^2} = \dfrac{y_p^2-y_0^2}{\left(a^2+\dfrac{y_p^2}{8}\right) - \left(a^2+\dfrac{y_0^2}{8}\right)} = 8$

            \item \ps 背景为极点极线，$M$在极线（准线$x = 1$）上，则$PQ$过极点($F_2(9, 0)$)，故可以先证明$PQ$过$F_2$，再利用双曲线的第一定义将$\bigtriangleup {F}_{1}{PQ}$ 周长转化到$|PQ|$上.

            根据$a = 3$可知双曲线方程为$\dfrac{x^2}{9} - \dfrac{y^2}{72} = 1$

            设$M(1, n)$，$n \in (-8\sqrt{2}, -4\sqrt{2}) \cup (4\sqrt{2}, 8\sqrt{2})$

            $A_1M: x = \dfrac{4y}{n} - 3$，$A_2M: x = -\dfrac{2y}{n}+3$

            联立$A_1M$和$\dfrac{x^2}{9} - \dfrac{y^2}{72} = 1$可得$\left(\dfrac{4y}{n}-3\right)^2 - \dfrac{y^2}{8} = 9$，一解为$y = 0$，另一解为$P$的纵坐标$y_p = \dfrac{192n}{128-n^2}$，代入直线方程可得$P\left(\dfrac{768}{128-n^2}-3, \dfrac{192n}{128-n^2}\right)$

            联立$A_2M$和$\dfrac{x^2}{9} - \dfrac{y^2}{72} = 1$可得$\left(3-\dfrac{2y}{n}\right)^2 - \dfrac{y^2}{8} = 9$，同理可得$Q\left(\dfrac{-192}{32-n^2}+3, \dfrac{96n}{32-n^2}\right)$

            故$\vv{F_2P} = \left(\dfrac{12n^2-768}{128-n^2}, \dfrac{192n}{128-n^2}\right), \vv{F_2Q} = \left(\dfrac{6n^2-384}{32-n^2}, \dfrac{96n}{32-n^2}\right)$

            故$\vv{F_2P} // \vv{F_2Q}$，故$P,F_2,Q$三点共线

            $C_{\bigtriangleup {F}_{1}{PQ}} = |PQ| + |PQ| + 12$

            设$PQ: x = my + 9$，和$\dfrac{x^2}{9} - \dfrac{y^2}{72} = 1$ 得 $\left( {8{m}^{2} - 1}\right) {y}^{2} + {144my} + 576 = 0$
            
            $\left\{  \begin{array}{l} 8{m}^{2} - 1 \neq  0 \\  \Delta  > 0 \end{array}\right.$，且$P,Q$均在右一支可得$m \in \left(-\dfrac{\sqrt{2}}{4}, \dfrac{\sqrt{2}}{4}\right)$
            
            由韦达定理可得 $\left\{\begin{array}{l}
                 {y}_{1} + {y}_{2} =  - \dfrac{144m}{8{m}^{2} - 1}  \\
                 {y}_{1}{y}_{2} = \dfrac{576}{8{m}^{2} - 1}
            \end{array}\right.$

            $|PQ| = \sqrt{1+m^2} \sqrt{(y_1+y_2)^2 - 4y_1y_2} = \sqrt{1 + {m}^{2}} \cdot  \dfrac{{48}\sqrt{1 + {m}^{2}}}{\left| 8{m}^{2} - 1\right| } = \dfrac{{48}\left( {1 + {m}^{2}}\right) }{1 - 8{m}^{2}} = \dfrac{-6\left( {1 - 8{m}^{2}}\right)  + {54}}{1 - 8{m}^{2}} = \dfrac{54}{1 - 8{m}^{2}} - 6 \in  \lbrack {48}, + \infty )$

            因此$C_{\bigtriangleup {F}_{1}{PQ}} = |PQ| + |PQ| + 12 \in [108, +\infty)$
        \end{enumerate}
    \end{jieda}
    \item 已知点 $B\left( {4,\sqrt{3}}\right)$ 是双曲线 $T : \dfrac{{x}^{2}}{{a}^{2}} - {y}^{2} = 1$ 上一点， $\Gamma$ 在点 $B$ 处的切线与 $x$ 轴交于点 $A$.
    \begin{enumerate}
        \item 求双曲线 $\Gamma$ 的方程及点 $A$ 的坐标
        \item 过 $A$ 且斜率非负的直线与 $\Gamma$ 的左、右支分别交于 $N,M$. 过 $N$ 做 ${NP}$ 垂直于 $x$ 轴交 $\Gamma$ 于 $P$ (当 $N$ 位于左顶点时认为 $N$ 与 $P$ 重合). $C$ 为圆 $E : {\left( x - 1\right) }^{2} + (y +  2)^{2} = 1$ 上任意一点，求四边形 ${MBPC}$ 的面积 $S$ 的最小值.
    \end{enumerate}
    \begin{jieda}
        \begin{enumerate}
            \item $a = 2$，$\dfrac{x^2}{4}-y^2 = 1$，$x - \sqrt{3}y - 1 = 0$，$A(1, 0)$
            \item \ps 背景为极点极线，做$M$关于$x$轴的对称点$M'$，则根据对称性可知$AP$过$M'$，根据极点极线的几何定义可知$MP, M'N$交点在$A$的极线$x = 4$上，根据对称性可知交点$Q$在$x$轴上，因此$PM$过定点$Q(4, 0)$
            \begin{center}
            \includegraphics[width=0.3\textwidth]{bodyxb/src/1_1089_1872_450_301_0.jpg}
            \end{center}
            先证明$PM$过定点$Q(4, 0)$：

            设$P(x_1, y_1), M(x_2, y_2)$，则$N(x_1, -y_1)$，当$k_{PM} \neq 0$时，设直线$PM: x = my + n$

            $\vv{AM} = (x_2-1, y_2) = (my_2+n-1, y_2)$，$\vv{AN} = (x_1-1, -y_1) = (my_1+n-1, -y_1)$

            $A, M, N$三点共线，因此有$(my_2+n-1)(-y_1) = (my_1+n-1)y_2$，即$2my_1y_2+(n-1)(y_1+y_2) = 0$
            
            联立直线和双曲线有$\left\{\begin{array}{l}
                x = my + n \\
                \dfrac{x^2}{4}-y^2 = 1
            \end{array}\right.$，故$(my+n)^2-4y^2=4$，即$(m^2-4)y^2+2mny+n^2-4 = 0$，满足$m^2 \neq 4, \Delta > 0$

            根据韦达定理有$\left\{\begin{array}{l}
                y_1+y_2 = \dfrac{-2mn}{m^2-4} \\
                y_1y_2 = \dfrac{n^2-4}{m^2-4}
            \end{array}\right.$，故$2m(n^2-4)+(n-1)(-2mn) = 0$，故$n = 4$

            $k_{PM} = 0$时，依然满足$PM$过定点$Q(4, 0)$，又因为$BQ \bot x$轴，所以 $S = S_{\triangle BMP} + S_{\triangle CMP} = S_{\triangle BQP} - S_{\triangle BQM} + S_{\triangle CMP} = \dfrac{1}{2}|BQ| \times |x_1 - x_2| + \dfrac{1}{2}|PM| \times h_{C-PM} \geq \dfrac{1}{2}|BQ| \times |x_1 - x_2| + \dfrac{1}{2}|PM| \times (h_{E-PM}-1)$

            易知当$PM$与$x$轴重合时，$|x_1-x_2|$最小，$|PM|$最小，$h_{E-PM}$最小，因此$S$此时最小，为$2+2\sqrt{3}$
        \end{enumerate}
    \end{jieda}

\end{enumerate}